%% Manuscript formatted for Applied Mathematical Finance (Taylor & Francis)
%% Requires the Taylor & Francis "interact" class file (interact.cls), which is
%% NOT bundled here: download it from the AMF "Instructions for authors ->
%% LaTeX templates" page and place interact.cls in this directory, then compile.
%% See AMF_submission.md for full instructions and the reference-style choice.
\documentclass[]{interact}

\usepackage{amsmath,amssymb,amsthm}
\usepackage{graphicx}
\usepackage{booktabs}
\usepackage{enumitem}
\usepackage[colorlinks=false,allbordercolors=white]{hyperref}
\usepackage{xurl}

\theoremstyle{plain}
\newtheorem{theorem}{Theorem}
\newtheorem{proposition}{Proposition}
\newtheorem{definition}{Definition}

\begin{document}

\articletype{Research Article}

\title{Swap Curve Construction and Optimization of the Tension Parameter:
An Energy-Partition Formulation}

\author{\name{Ashutosh Sinha\thanks{Email: ajsinha@gmail.com}}
\affil{Independent Researcher}}

\maketitle

\begin{abstract}
\noindent Polynomial splines have long been used for the construction of swap curves. They
suffer, however, from excessive convexity, a lack of locality under perturbation of the
input rates, and the possibility of negative implied forwards. Tension splines remedy these
defects through a single shape parameter, but the choice of that parameter has remained
largely heuristic. In this paper we place the tension spline on its proper physical footing
as the equilibrium shape of a thin elastic beam held under longitudinal tension, and we
derive in closed form the two energies stored by such a beam: a \emph{flexural} (bending)
energy and an \emph{axial} (tension) energy. We show that the curvature ``smoothness'' norm
used in the literature \emph{is} the flexural energy, and that the first-derivative norm
\emph{is} the axial energy, so that the two are the complementary halves of one elastic
energy rather than independent quantities. We then propose selecting the tension parameter
by the \emph{bending-energy fraction}
$\varphi(\sigma)=E_{\mathrm{bend}}/(E_{\mathrm{bend}}+E_{\mathrm{tens}})$, a dimensionless
quantity that runs monotonically from one (cubic) to zero (linear). The neutral choice
$\varphi(\sigma^\star)=\tfrac12$ corresponds to equipartition of elastic energy; the target
$\varphi^\star$ provides a single, transparent dial through which a desk may express its
trade-off between smoothness and locality. The criterion is invariant under rescaling of the
rate and time axes, removing the unit dependence of earlier energy-based schemes. On the U.S.
Treasury par-yield curve of $22$-May-$2026$ the criterion yields a unique optimum
($\sigma^\star=4.86$ at equipartition) and, relative to the natural cubic spline, reduces the
hedge-locality radius of the benchmark $5$-year node from $6.7$ years to $0.8$ year --- a
factor of eight --- while suppressing the spurious long-end forward hump that the cubic
manufactures from a flat $20$--$30$-year segment. Across $1{,}064$ weekly U.S. Treasury curves
spanning 2006--2026, the selected tension tracks the rate regime smoothly without jumps, and the
predicted inverse relationship between tension and hedge-locality radius holds out of sample with
a log-correlation of $-0.99$.
\end{abstract}

\begin{keywords}
yield curve construction; term-structure interpolation; tension spline;
energy-partition criterion; hedge locality; instantaneous forward curve;
multi-curve SOFR/OIS; variational method; Euler--Bernoulli beam
\end{keywords}

\begin{jelcode}
C63; G12; E43; C61; G17
\end{jelcode}

\section{Introduction}
A risk-free term structure is fundamental to the pricing of fixed-income securities and to
the risk management of the same. The term structure is required to be a continuous function
of time, with additional requirements on the continuity of its first and second derivatives.
Given these constraints, computational procedures are needed to estimate the structure, and
for those procedures the Treasury rates alone are not enough: to reflect the market a
combination of instruments --- cash deposits, futures, swap rates --- is used to estimate the
curve.

Finding the right set of instruments is only half the problem. As Hagan and West~\cite{Hagan}
articulate, there is no single best way to construct a term structure from a set of rates.
It is commonly desired that the derived curve be smooth, that the interpolation be
\emph{local}, that the implied forwards be both smooth and stable, and that hedges be as
local as possible. Smoothness means the absence of kinks and spurious wiggles; market
fidelity means the curve reproduces known market rates exactly; locality means that a
perturbation of one input propagates to a bounded neighbourhood of the curve, which is
precisely what permits a fixed-income portfolio to be hedged with a small, well-localised set
of instruments.

Parametric methods such as Nelson--Siegel produce a good overall shape but leave much to be
desired in the accuracy of fit. Cubic regression splines, first applied to the term structure
by McCulloch~\cite{McCulloch}, give the smoothest interpolant and have been shown to be at
least as satisfactory as competing methods. Cubic splines are not, however, free of defects:
they fail to guarantee monotonicity or convexity of the data~\cite{Anderson}, they are
non-local, and for rate curves they can produce negative forwards. Tanggaard~\cite{Tanggaard}
introduced non-parametric smoothing of the yield curve through a penalised least-squares
estimator; Anderson~\cite{Anderson} extended this idea, modifying the penalty term to obtain
hyperbolic tension splines. Kamada and Enkhbat~\cite{Kamada} present the tension spline as a
generalisation that interpolates between the cubic (minimiser of the squared second
derivative) and the linear (minimiser of the squared first derivative) splines, and so
provides a convenient middle ground that preserves shape~\cite{Costantini}. Costantini, Kvasov
and Manni~\cite{Costantini} treat the hyperbolic tension spline as the solution of a
differential multipoint boundary-value problem; Renka~\cite{Renka} automates the choice of the
tension factor subject to monotonicity and convexity constraints.

The practical relevance of these defects is underscored by official practice. The U.S.
Treasury, which publishes the most widely referenced government curve in the world, abandoned
its quasi-cubic Hermite spline in December 2021 and now estimates its daily par-yield curve
with a monotone convex spline~\cite{TreasuryMethod}, precisely to suppress the oscillation and
non-locality of the cubic. The tension spline occupies exactly the continuum between these two
regimes: it relaxes to the cubic at zero tension and tightens towards the shape-preserving,
local, piecewise-linear limit as the tension grows, with a single parameter governing the
position along that continuum.

We make two methodological observations about this body of work. First, automated selection
of the tension factor, while elegant, does not by itself respond to a practitioner's risk and
hedging preferences; any defensible selection rule should expose those preferences as an
explicit parameter rather than hide them inside a fixed constraint set. Second --- and this is
the technical heart of the paper --- the ``elastic beam'' analogy that motivates the tension
spline can be made fully rigorous. The tension spline is not merely \emph{like} a beam; it
\emph{is} the equilibrium shape of a thin elastic beam under longitudinal tension, and the
norms that the literature uses to score a curve are the literal flexural and axial energies of
that beam. Recognising this collapses the two seemingly independent ``smoothness'' and
``length'' penalties into the two halves of a single stored energy, and it suggests a natural,
dimensionless, range-independent rule for selecting the tension. This paper revises and
substantially extends an earlier treatment by the author, recasting that elastic-beam heuristic
on a rigorous variational footing and replacing the ad hoc selection rule with the
dimensionless, scale-invariant criterion developed here.

A note on scope. For clarity the core construction is presented on a single curve; current
market practice discounts on an overnight (OIS/SOFR) curve within a multi-curve framework
\cite{Bianchetti,Henrard}, and Section~\ref{sec:multicurve} demonstrates the method directly in
that setting. The interpolation theory developed below is agnostic to that choice: it acts
downstream of the bootstrap and applies unchanged to any single curve, whether a projection
curve or a discount curve. We also interpolate in the space of instantaneous forward rates, following the
emphasis of Hagan and West~\cite{Hagan} on the importance of the interpolation variable; the
forward space is the most demanding test of an interpolant because forwards expose, through the
derivative, the oscillations that zero rates conceal.

The remainder of the paper is organised as follows. Section~\ref{sec:fi} reviews fixed-income
pricing in relation to vanilla swaps. Section~\ref{sec:cubic} reviews the natural cubic spline.
Section~\ref{sec:tension} establishes the tension spline as a tensioned elastic beam and
derives its variational characterisation. Section~\ref{sec:energy} derives the flexural and
axial energies in closed form and presents the energy-partition selection criterion.
Section~\ref{sec:validation} validates the criterion through forward positivity and hedge
locality. Section~\ref{sec:fwd} extends the method to the instantaneous forward curve,
Section~\ref{sec:nonunif} to per-segment non-uniform tension, and Section~\ref{sec:fwdnu}
unifies the two. Section~\ref{sec:multicurve} embeds the method in a multi-curve SOFR/OIS
construction. Section~\ref{sec:risk} derives the analytic risk Jacobians and the stability of the
dial. Section~\ref{sec:empirical} benchmarks the method against production interpolators and
characterises $\varphi^\star$. Section~\ref{sec:why} sets out the
practical benefits, and Section~\ref{sec:discussion} discusses results and future directions.
Appendix~\ref{app:nonunif} derives the non-uniform spline, Appendix~\ref{app:energy} the
closed-form energy integrals, and Appendix~\ref{app:stable} the numerically stable evaluation.

\section{Fixed-Income Pricing Basics}\label{sec:fi}
A money-market account evolves as $dB(t)=r_t\,B(t)\,dt$, so that
$B(t)=\exp\!\big(\int_0^t r_s\,ds\big)$ and the stochastic discount factor is
$D(t,T)=\exp\!\big(-\int_t^T r_s\,ds\big)$, where $r_t$ is the instantaneous short rate. The
zero-coupon bond price $P(t,T)$ is the risk-neutral expectation of $D(t,T)$ and equals
$D(t,T)$ when the rate process is deterministic. With $\tau(t,T)$ the year fraction to
maturity, the continuously compounded zero rate $R(t,T)$ satisfies
\begin{equation}
P(t,T)=e^{-R(t,T)\,\tau(t,T)},\qquad
R(t,T)=-\frac{\ln P(t,T)}{\tau(t,T)}.
\end{equation}
Simply-compounded market rates (for example term SOFR) accrue on an Act/360 basis,
\begin{equation}
L(t,T)=\frac{1-P(t,T)}{\tau(t,T)\,P(t,T)} .
\end{equation}
The continuously compounded forward rate for $t\le T_1\le T_2$ is
\begin{equation}
F(t,T_1,T_2)=\frac{1}{\tau(T_1,T_2)}\ln\frac{P(t,T_1)}{P(t,T_2)},
\end{equation}
and for simple compounding
\begin{equation}
F(t,T_1,T_2)=\frac{1}{\tau(T_1,T_2)}\!\left(\frac{P(t,T_1)}{P(t,T_2)}-1\right). \label{eq:fwdsimple}
\end{equation}
As $T_2\to T_1$ the forward rate becomes the instantaneous forward
$f(t,T)=-\partial_T\ln P(t,T)$, whence $P(t,T)=\exp\!\big(-\int_t^T f(t,u)\,du\big)$ and
$F(t,T_1,T_2)=\tfrac{1}{T_2-T_1}\int_{T_1}^{T_2} f(t,u)\,du$.

A vanilla payer swap exchanges a fixed leg for a floating leg from a future start date. Its
defining attributes --- notional (used only for cash-flow calculation), fixed rate, coupon
frequency, business-day convention, floating index, day-count, effective and maturity dates ---
are standard. A swap has zero present value at inception; the swap rate is the fixed rate that
makes it so. Bootstrapping the curve from cash, futures and swaps proceeds, following
Hagan~\cite{Hagan}, through the recursion for the discount factor implied by the $n$-th swap,
\begin{equation}
B(t,t_n)=\frac{1-R_n\sum_{j=1}^{n-1}\theta_j B(t,t_j)}{1+R_n\theta_n},\qquad
r_n=-\frac{1}{\tau_n}\ln B(t,t_n),
\label{eq:bootstrap}
\end{equation}
where $R_n$ is the par swap rate, $\theta_j$ the coupon year fractions, and the unit in the
numerator should strictly be the discount factor to the (two-business-day) effective date,
here approximated by one. Short-end deposit and futures stubs are converted to continuously
compounded spot rates through $f_{sl}=\theta_{sl}^{-1}\ln\big[(1+r_l\theta_l)/(1+r_s\theta_s)\big]$
and $(1+F\theta_F)=e^{f\theta_f}$ respectively.

Applying \eqref{eq:bootstrap}, or in the illustrative case below simply reading the published
constant-maturity yields, produces the curve nodes used throughout this paper.

For a current, fully reproducible illustration we take the U.S. Treasury par-yield curve of
$22$-May-$2026$ (Table~\ref{tab:data}), as published in the Federal Reserve H.15
release~\cite{H15}. The eleven constant-maturity points span one month to thirty years and
carry two features that stress any interpolant: a front-end kink, the one-month yield
($3.72\%$) sitting above the three-month ($3.68\%$) before the curve resumes rising, and a
nearly flat long end, the twenty- and thirty-year yields almost coinciding at $5.06$--$5.07\%$.
We interpolate the observed yield curve directly; treating the par yields as a zero curve for
the purpose of deriving instantaneous forwards is an approximation made only for the diagnostic
of Section~\ref{sec:validation}, and the interpolation theory is unchanged if one instead
tensions a bootstrapped zero or forward curve.

\begin{table}[t]\small\centering
\caption{U.S. Treasury CMT par-yield curve, $22$-May-$2026$ (Federal Reserve H.15).}\label{tab:data}
\begin{tabular}{lrrrrrrrrrrr}
\toprule
Tenor & 1M & 3M & 6M & 1Y & 2Y & 3Y & 5Y & 7Y & 10Y & 20Y & 30Y\\
\midrule
Yield (\%) & 3.72 & 3.68 & 3.79 & 3.86 & 4.13 & 4.18 & 4.27 & 4.41 & 4.56 & 5.06 & 5.07\\
\bottomrule
\end{tabular}
\end{table}

\section{The Natural Cubic Spline}\label{sec:cubic}
For knots $x_0<x_1<\dots<x_{n-1}$ with data $y_i$, the natural cubic spline is the function
that is cubic on each $[x_i,x_{i+1}]$, interpolates the data, and has continuous first and
second derivatives. Writing $z_i=f''(x_i)$ and $h_i=x_{i+1}-x_i$, the requirement
$f^{(\mathrm{iv})}=0$ on each interval makes $f''$ linear; the linear function taking the values
$z_i$ at $x_i$ and $z_{i+1}$ at $x_{i+1}$ is
\begin{equation}
f_i''(x)=\frac{z_i}{h_i}(x_{i+1}-x)+\frac{z_{i+1}}{h_i}(x-x_i).
\end{equation}
Integrate once, then again, carrying two constants $p_i,q_i$:
\[
f_i'(x)=-\frac{z_i}{2h_i}(x_{i+1}-x)^2+\frac{z_{i+1}}{2h_i}(x-x_i)^2+p_i,\qquad
f_i(x)=\frac{z_i}{6h_i}(x_{i+1}-x)^3+\frac{z_{i+1}}{6h_i}(x-x_i)^3+p_i(x-x_i)+q_i.
\]
Imposing $f_i(x_i)=y_i$ gives $q_i=y_i-z_ih_i^2/6$, and $f_i(x_{i+1})=y_{i+1}$ gives
$p_i=(y_{i+1}-y_i)/h_i-h_i(z_{i+1}-z_i)/6$. Collecting powers of $(x-x_i)$ yields the Horner form
\begin{equation}
f_i(x)=y_i+(x-x_i)\big[C_i+(x-x_i)\big(B_i+(x-x_i)A_i\big)\big],
\end{equation}
\begin{equation}
A_i=\frac{z_{i+1}-z_i}{6h_i},\quad B_i=\frac{z_i}{2},\quad
C_i=\frac{y_{i+1}-y_i}{h_i}-\frac{h_i}{6}z_{i+1}-\frac{h_i}{3}z_i.
\end{equation}
By construction $f$ and $f''$ are already continuous at the interior knots (both sides share
$y_i$ and $z_i$). The remaining condition is continuity of $f'$ at each $x_i$. From the integrated
slope, the left interval $i-1$ gives
$f_{i-1}'(x_i)=p_{i-1}+z_ih_{i-1}/2=(y_i-y_{i-1})/h_{i-1}+h_{i-1}(2z_i+z_{i-1})/6$, and the right
interval $i$ gives $f_i'(x_i)=p_i-z_ih_i/2=(y_{i+1}-y_i)/h_i-h_i(2z_i+z_{i+1})/6$. Equating
$f_{i-1}'(x_i)=f_i'(x_i)$ and multiplying by $6$ gives the diagonally dominant tridiagonal system,
for $i=1,\dots,n-2$,
\begin{equation}
h_{i-1}z_{i-1}+2(h_{i-1}+h_i)z_i+h_i z_{i+1}
=6\!\left(\frac{y_{i+1}-y_i}{h_i}-\frac{y_i-y_{i-1}}{h_{i-1}}\right),
\end{equation}
with the natural boundary conditions $z_0=z_{n-1}=0$. The cubic spline minimises
$\int_a^b\!\big[f''(x)\big]^2\,dx$; since the curvature of a graph is
$|f''|\,[1+(f')^2]^{-3/2}$ and $(f')^2$ is negligible for a rate curve, the cubic is
effectively the minimum-curvature interpolant. Its smoothness is, however, bought at the cost
of locality and --- for steep or kinked data --- of possible overshoot.

\section{The Tension Spline as a Tensioned Elastic Beam}\label{sec:tension}
\subsection{The physical model}
The name ``spline'' is borrowed from the draughtsman's tool: a thin elastic lath threaded
through fixed pins. Left free, the lath assumes the minimum-bending-energy shape --- the cubic
spline. If the lath is additionally pulled taut by a longitudinal force $T$ at its ends, it
straightens between the pins; in the limit of infinite tension it becomes the polygonal
(linear) interpolant. This is precisely the behaviour wanted of a rate curve: a tunable
departure from the over-convex cubic towards the economically conservative piecewise-linear
shape. We therefore model the curve as a thin elastic beam of flexural rigidity $EI$ held under
axial tension $T$, and we make the analogy exact rather than illustrative.

Within the small-deflection (Euler--Bernoulli) theory, the elastic energy stored by such a beam
over $[a,b]$ is the sum of a flexural term and an axial (tension) term,
\begin{equation}
\mathcal{E}[f]=\int_a^b\!\left[\tfrac12 EI\,\big(f''\big)^2+\tfrac12 T\,\big(f'\big)^2\right]dx .
\label{eq:functional}
\end{equation}
The first term is the bending strain energy $\int M^2/(2EI)\,dx$ with bending moment
$M=EI\,f''$; the second is the work done against the axial force as the beam elongates,
$T$ times the second-order arc-length stretch $\tfrac12\int (f')^2\,dx$. To see the latter,
the arc length of the graph is $\int\sqrt{1+(f')^2}\,dx=\int\big[1+\tfrac12(f')^2+O((f')^4)\big]dx$,
so the elongation beyond the chord is $\tfrac12\int(f')^2\,dx$ to leading order, and the work
done by the constant tension $T$ against this elongation is $\tfrac12 T\int(f')^2\,dx$. The two
energies, the moments and the deflection are shown in Figure~\ref{fig:fbd}.

For clarity, the symbols carry the following meaning. The function $f(x)$ is the interpolated
rate curve itself: the height of the beam above the horizontal axis at abscissa $x$ (a tenor),
i.e.\ the rate at that tenor; its first derivative $f'(x)$ is the slope of the curve and its
second derivative $f''(x)$ is the curvature. The constant $E$ is the material's Young's modulus
(stiffness of the beam material) and $I$ is the second moment of area of its cross-section; their
product $EI$, the \emph{flexural rigidity}, measures how strongly the beam resists bending --- a
large $EI$ yields a smooth, gently curved shape. The constant $T$ is the axial \emph{tension},
the longitudinal force pulling the beam taut at its ends; a large $T$ straightens the beam
between the knots. The bending moment $M(x)=EI\,f''(x)$ is the internal torque the beam carries
at $x$. Only the dimensionless combination of the two physical constants matters for the curve
shape: the \emph{tension parameter}
\[
\sigma=\sqrt{T/EI}\quad(\text{units of }1/x),
\]
which sets the single knob of the construction. When $\sigma=0$ (no tension, $T=0$) the
functional \eqref{eq:functional} retains only the bending term and $f$ is the natural cubic
spline; as $\sigma\to\infty$ (overwhelming tension) the axial term dominates and $f$ approaches
the piecewise-linear interpolant. The limits $a,b$ are the first and last tenors, and the
$\tfrac12$ factors are the usual quadratic-energy convention. Because only $\sigma$ enters the
Euler--Lagrange equation below, the individual values of $E$, $I$ and $T$ never need to be
specified --- they are absorbed into the one parameter $\sigma$.

\begin{figure}[h]\centering
\includegraphics[width=0.7\textwidth]{fig_fbd.png}
\caption{The tension-spline element as a thin elastic beam under axial tension $T$. The bending
moment $M=EI f''$ stores the flexural energy $\tfrac12 EI\!\int(f'')^2$ (the ``smoothness''
norm); the slope $f'$ stores the axial energy $\tfrac12 T\!\int(f')^2$ (the ``length'' norm).
These are the two halves of one stored elastic energy.}\label{fig:fbd}
\end{figure}

\begin{proposition}\label{prop:euler}
The stationary curve of \eqref{eq:functional} subject to interpolation
$f(x_i)=y_i$ satisfies, on each open interval $(x_i,x_{i+1})$,
\begin{equation}
EI\,f^{(\mathrm{iv})}-T\,f''=0
\qquad\Longleftrightarrow\qquad
f^{(\mathrm{iv})}-\sigma^2 f''=0,\qquad \sigma^2:=\frac{T}{EI}.
\label{eq:euler}
\end{equation}
\end{proposition}
\begin{proof}
Perturb $f\mapsto f+\epsilon\,\delta f$ with $\delta f(x_i)=0$ at every knot (interpolation is
fixed). Differentiating \eqref{eq:functional} and evaluating at $\epsilon=0$,
\[
\delta\mathcal{E}=\int_a^b\!\big[EI\,f''\,\delta f''+T\,f'\,\delta f'\big]\,dx .
\]
Integrate the flexural term by parts twice:
\[
\int_a^b EI\,f''\,\delta f''\,dx
=\big[EI\,f''\,\delta f'\big]_a^b-\int_a^b EI\,f'''\,\delta f'\,dx
=\big[EI\,f''\,\delta f'-EI\,f'''\,\delta f\big]_a^b+\int_a^b EI\,f^{(\mathrm{iv})}\,\delta f\,dx,
\]
and the axial term by parts once:
$\int_a^b T f'\,\delta f'\,dx=\big[T f'\,\delta f\big]_a^b-\int_a^b T f''\,\delta f\,dx$.
The boundary terms in $\delta f$ vanish because $\delta f=0$ at the knots; the remaining
boundary terms $\big[EI f''\,\delta f'\big]$ vanish under the natural free-end condition
$f''(a)=f''(b)=0$. Collecting the interior integrals,
\[
\delta\mathcal{E}=\int_a^b\!\big[EI\,f^{(\mathrm{iv})}-T\,f''\big]\,\delta f\,dx .
\]
Since this must vanish for every admissible $\delta f$, the fundamental lemma of the calculus of
variations gives $EI f^{(\mathrm{iv})}-T f''=0$ on each open interval, i.e.\ \eqref{eq:euler}.
\end{proof}

Equation~\eqref{eq:euler} is the hyperbolic tension-spline equation of Schweikert~\cite{Schweikert}
and Cline~\cite{Cline}; the tension parameter $\sigma$ that the literature introduces by fiat is
here identified as $\sigma=\sqrt{T/EI}$, the square root of the dimensionless ratio of axial to
flexural stiffness. As $\sigma\to0$ (no tension) the beam relaxes to the cubic; as
$\sigma\to\infty$ it pulls taut to the linear interpolant.

\subsection{Construction}
The characteristic polynomial of \eqref{eq:euler} is $r^4-\sigma^2 r^2=r^2(r^2-\sigma^2)=0$,
with roots $0,0,\pm\sigma$, so the general solution on $[x_i,x_{i+1}]$ is
\[
f(x)=a_i+b_i x+c_i\sinh(\sigma x)+d_i\cosh(\sigma x).
\]
Differentiating twice, $f''=\sigma^2\big(c_i\sinh\sigma x+d_i\cosh\sigma x\big)$, so $f''$ is a
pure hyperbolic combination; writing $z_i:=f''(x_i)$ and $z_{i+1}:=f''(x_{i+1})$ and solving the
$2\times2$ system for the hyperbolic part gives
$f''(x)=\big[z_i\sinh(\sigma(x_{i+1}-x))+z_{i+1}\sinh(\sigma(x-x_i))\big]/\sinh(\sigma h_i)$,
which is \eqref{eq:t2deriv} and already satisfies $f''(x_i)=z_i$, $f''(x_{i+1})=z_{i+1}$.
Integrating $f''$ twice and fixing the two constants of integration by the interpolation
conditions $f(x_i)=y_i$, $f(x_{i+1})=y_{i+1}$ removes the homogeneous linear part in favour of
the chord, giving \eqref{eq:tspline}; one differentiation gives \eqref{eq:tderiv}. Explicitly,
with $h_i=x_{i+1}-x_i$,
\begin{align}
f(x)&=\frac{z_i\sinh\!\big(\sigma(x_{i+1}-x)\big)+z_{i+1}\sinh\!\big(\sigma(x-x_i)\big)}
{\sigma^2\sinh(\sigma h_i)}
+\Big(y_i-\tfrac{z_i}{\sigma^2}\Big)\frac{x_{i+1}-x}{h_i}
+\Big(y_{i+1}-\tfrac{z_{i+1}}{\sigma^2}\Big)\frac{x-x_i}{h_i},\label{eq:tspline}\\[2pt]
f'(x)&=\frac{-z_i\cosh\!\big(\sigma(x_{i+1}-x)\big)+z_{i+1}\cosh\!\big(\sigma(x-x_i)\big)}
{\sigma\sinh(\sigma h_i)}
+\frac{(y_{i+1}-y_i)-(z_{i+1}-z_i)/\sigma^2}{h_i},\label{eq:tderiv}\\[2pt]
f''(x)&=\frac{z_i\sinh\!\big(\sigma(x_{i+1}-x)\big)+z_{i+1}\sinh\!\big(\sigma(x-x_i)\big)}
{\sinh(\sigma h_i)}.\label{eq:t2deriv}
\end{align}
The construction makes $f$ and $f''$ continuous at the interior knots automatically (both sides
share $y_i$ and $z_i$); the remaining $n-2$ unknowns $z_1,\dots,z_{n-2}$ are fixed by requiring
continuity of $f'$. Evaluating \eqref{eq:tderiv} at $x_i^{-}$ (from interval $i-1$) and $x_i^{+}$
(from interval $i$) and equating yields, after collecting the coefficients of $z_{i-1},z_i,z_{i+1}$,
the diagonally dominant tridiagonal system, for $i=1,\dots,n-2$,
\begin{equation}
\alpha_{i-1}z_{i-1}+(\beta_{i-1}+\beta_i)z_i+\alpha_i z_{i+1}=\gamma_i-\gamma_{i-1},
\label{eq:ttridiag}
\end{equation}
\begin{equation}
\alpha_i=\frac{1}{h_i}-\frac{\sigma}{\sinh(\sigma h_i)},\quad
\beta_i=\frac{\sigma\cosh(\sigma h_i)}{\sinh(\sigma h_i)}-\frac{1}{h_i},\quad
\gamma_i=\frac{\sigma^2(y_{i+1}-y_i)}{h_i},
\end{equation}
with $z_0=z_{n-1}=0$. As $\sigma\to0$, Taylor expansion gives $\alpha_i\to -h_i/6$,
$\beta_i\to h_i/3$ and $\gamma_i\to (y_{i+1}-y_i)/h_i$ (times $\sigma^2$, which cancels), so
\eqref{eq:ttridiag} reduces to the natural-cubic system of Section~\ref{sec:cubic}; numerically
at $\sigma=10^{-3}$ the two interpolants of Table~\ref{tab:data} agree to $3\times10^{-9}$. The
implementation differs from the cubic only in the hyperbolic evaluations; for large $\sigma h_i$
the formulae are reorganised into the stable $e^{-\sigma h}$ form of Appendix~\ref{app:stable} to
avoid overflow. Figure~\ref{fig:curves} shows the linear, cubic, and two
tension interpolants of the yield curve; smaller $\sigma$ lies near the cubic and larger
$\sigma$ approaches the linear shape. Already in yield space the cubic overshoots: it lifts the
flat $20$--$30$-year segment into a spurious hump of roughly five basis points peaking near
twenty-four years, a feature absent from the data, whereas the tension curve hugs the flat
segment.

\begin{figure}[t]\centering
\includegraphics[width=0.78\textwidth]{fig_curves.png}
\caption{Linear, cubic, and tension-spline interpolation of the U.S. Treasury par-yield curve
of Table~\ref{tab:data}. The tension curve at the equipartition optimum $\sigma^\star=4.86$
(Section~\ref{sec:energy}) removes the cubic's spurious long-end hump while retaining a smooth
belly.}\label{fig:curves}
\end{figure}

\section{Energy Decomposition and Tension Selection}\label{sec:energy}
\subsection{The two energies are the two norms}
The literature scores a candidate curve by a smoothness norm $\int(f'')^2$ and, in tension
schemes, by a first-derivative norm weighted by $\sigma^2$. Proposition~\ref{prop:euler}
identifies these immediately. Define, per unit flexural rigidity, the \emph{flexural} and
\emph{axial} energies of the tensioned beam,
\begin{equation}
E_{\mathrm{bend}}(\sigma)=\tfrac12\!\int_0^T\!\big(f''_\sigma\big)^2\,dt,
\qquad
E_{\mathrm{tens}}(\sigma)=\tfrac12\,\sigma^2\!\int_0^T\!\big(f'_\sigma\big)^2\,dt .
\label{eq:energies}
\end{equation}
The smoothness norm \emph{is} $2E_{\mathrm{bend}}$ and the $\sigma^2$-weighted length norm
\emph{is} $2E_{\mathrm{tens}}$; their sum is the total stored elastic energy
$\mathcal{E}[f_\sigma]$ of \eqref{eq:functional}. They are therefore not independent scores to
be balanced against each other by an external rule, but the two complementary components of one
energy. This observation also dissolves a circularity in earlier ``beam-energy'' arguments that
compared a curvature norm against a separately constructed point-load deflection energy: by
Clapeyron's theorem the work done by a load equals the strain energy it stores, so the two were
estimates of the same flexural quantity, and the point-load route --- which requires inverting
for the location and magnitude of a fictitious load --- is both redundant and ill-posed for the
S-shaped segments a spline can produce. We discard it in favour of \eqref{eq:energies}.

Both integrals in \eqref{eq:energies} are elementary. Using \eqref{eq:t2deriv} and the
identities $\int_0^h\sinh^2(\sigma u)\,du=\tfrac{\sinh(2\sigma h)}{4\sigma}-\tfrac h2$ and
$\int_0^h\sinh(\sigma(h-u))\sinh(\sigma u)\,du=\tfrac12\big(h\cosh(\sigma h)-\sigma^{-1}\sinh(\sigma h)\big)$,
the per-segment flexural energy is
\begin{equation}
E^{(i)}_{\mathrm{bend}}=\frac{1}{2\sinh^2(\sigma h_i)}
\left[(z_i^2+z_{i+1}^2)\!\left(\frac{\sinh(2\sigma h_i)}{4\sigma}-\frac{h_i}{2}\right)
+z_i z_{i+1}\!\left(h_i\cosh(\sigma h_i)-\frac{\sinh(\sigma h_i)}{\sigma}\right)\right].
\label{eq:Eb}
\end{equation}
The axial energy admits the companion closed form \eqref{eq:Ea} of Appendix~\ref{app:energy}. Both expressions were verified against high-order quadrature to a
relative error of $3\times10^{-7}$. The total energies are
$E_{\mathrm{bend}}=\sum_i E^{(i)}_{\mathrm{bend}}$ and likewise for $E_{\mathrm{tens}}$.

\subsection{The selection criterion}
Every value of $\sigma$ already minimises its own functional \eqref{eq:functional}; the optimal
tension is therefore a question \emph{outside} the variational principle, and a selection rule
must supply an external preference. We make that preference dimensionless and explicit through
the \emph{bending-energy fraction}
\begin{equation}
\varphi(\sigma)=\frac{E_{\mathrm{bend}}(\sigma)}{E_{\mathrm{bend}}(\sigma)+E_{\mathrm{tens}}(\sigma)}\in(0,1).
\label{eq:phi}
\end{equation}
At $\sigma\to0$ the axial energy vanishes and $\varphi\to1$: the stored energy is pure
curvature, the cubic limit. At $\sigma\to\infty$ the interior curvature vanishes and
$\varphi\to0$: the energy is pure tension, the linear limit. The following proposition makes
the well-posedness of the selection precise.

\begin{proposition}[Properties of $\varphi$]\label{prop:phi}
Write $t=\sigma^2$, $B(\sigma)=\int(f_\sigma'')^2$, $A(\sigma)=\int(f_\sigma')^2$, so
$E_{\mathrm{bend}}=\tfrac12 B$, $E_{\mathrm{tens}}=\tfrac12 t A$ and
$\varphi(\sigma)=B/(B+tA)$. Then
\begin{enumerate}[leftmargin=2em,itemsep=1pt]
\item[(i)] $\varphi$ is continuous on $(0,\infty)$, with $\lim_{\sigma\to0^+}\varphi=1$ and
$\lim_{\sigma\to\infty}\varphi=0$;
\item[(ii)] consequently, for every target $\varphi^\star\in(0,1)$ there exists
$\sigma^\star>0$ with $\varphi(\sigma^\star)=\varphi^\star$;
\item[(iii)] the minimum stored energy $g(t)=\min_f\{B[f]+tA[f]\}=B(\sigma)+tA(\sigma)$ is
concave in $t$, and $A(\sigma)=g'(t)$ is non-increasing in $\sigma$;
\item[(iv)] $\varphi$ is strictly decreasing wherever $A\,B+t\,g''(t)\,g(t)>0$; this holds
throughout the range of interest, so $\sigma^\star$ is unique.
\end{enumerate}
\end{proposition}
\begin{proof}
\emph{(i)} The map $\sigma\mapsto f_\sigma$ is continuous (the tridiagonal system
\eqref{eq:ttridiag} has coefficients analytic in $\sigma$ and is uniformly invertible), so
$B,A$ are continuous and $\varphi$ is continuous where $B+tA>0$, i.e.\ for all $\sigma>0$. As
$\sigma\to0^+$, $f_\sigma$ tends to the natural cubic with $0<B\to B_{\mathrm{cub}}<\infty$ and
$A\to A_{\mathrm{cub}}<\infty$, so $tA\to0$ and $\varphi\to1$. As $\sigma\to\infty$, $f_\sigma$
tends to the piecewise-linear interpolant: $A\to\int(f_{\mathrm{lin}}')^2>0$ is bounded below,
while the curvature concentrates in boundary layers of width $O(1/\sigma)$ about the knots where
$f''=O(\sigma)$, so $B=O(\sigma)$. Hence $tA/B=O(\sigma^2)/O(\sigma)\to\infty$ and
$\varphi=1/(1+tA/B)\to0$.
\emph{(ii)} Immediate from (i) and the intermediate value theorem.
\emph{(iii)} For each fixed admissible $f$, $t\mapsto B[f]+tA[f]$ is affine with slope
$A[f]\ge0$; $g$ is the pointwise minimum of this family and is therefore concave. By the
envelope theorem $g'(t)=A[f_t]=A(\sigma)$, and concavity gives $g''=A'\le0$, so $A$ is
non-increasing.
\emph{(iv)} Let $R=tA/B=E_{\mathrm{tens}}/E_{\mathrm{bend}}$, so $\varphi=1/(1+R)$ and $\varphi$
decreases iff $R$ increases. Writing $A=g'$ and $B=g-tg'$,
\[
R'(t)=\frac{d}{dt}\frac{tg'}{g-tg'}=\frac{g'(g-tg')+t\,g''g}{(g-tg')^2}=\frac{A\,B+t\,g''g}{B^2},
\]
which is positive precisely under the stated condition. Since $g''\le0$ the condition can fail
only if the trade-off curve is strongly concave; on the curve of Table~\ref{tab:data} we verify
$d\varphi/d\sigma<0$ at every point of a dense grid, and (i) guarantees the endpoints, so
$\sigma^\star$ is unique.
\end{proof}

We state plainly what is and is not proved: parts (i)--(iii) hold for any non-degenerate data;
the strict global monotonicity (iv) is reduced to an explicit sign condition that we verify
numerically rather than assert as an unconditional theorem, since for adversarial knot
configurations the trade-off curve $g$ could in principle be concave enough to violate it.
Existence of $\sigma^\star$ is unconditional; uniqueness is guaranteed under (iv).

\begin{definition}[Energy-partition tension]
Given a target fraction $\varphi^\star\in(0,1)$, the optimal tension $\sigma^\star$ is the
root of $\varphi(\sigma)=\varphi^\star$. The neutral choice $\varphi^\star=\tfrac12$
corresponds to \emph{equipartition} of elastic energy between flexure and tension.
\end{definition}

Two further properties recommend \eqref{eq:phi} over earlier max-normalised rules. First, it is
\emph{range-independent}: $\varphi$ is an intrinsic property of each $\sigma$ and does not
depend on the window of $\sigma$ over which a normalisation is taken. Second, it is
\emph{scale-invariant}.

\begin{proposition}[Scale invariance]\label{prop:inv}
$\varphi(\sigma)$ is invariant under rescaling of the rate axis, and under joint rescaling of
the time axis with $\sigma$ carried in inverse-time units.
\end{proposition}
\begin{proof}
\emph{Rate axis.} Replace $y_i\mapsto\alpha y_i$. The tridiagonal right-hand side
$\gamma_i-\gamma_{i-1}$ in \eqref{eq:ttridiag} scales by $\alpha$ while the matrix is unchanged,
so $z_i\mapsto\alpha z_i$ and hence $f\mapsto\alpha f$, $f'\mapsto\alpha f'$, $f''\mapsto\alpha f''$.
Then $B=\int(f'')^2\mapsto\alpha^2 B$ and $A=\int(f')^2\mapsto\alpha^2 A$, so both energies scale
by $\alpha^2$ and $\varphi=B/(B+\sigma^2A)$ is unchanged. This is the basis-points-versus-decimals
invariance.
\emph{Time axis.} Replace $x\mapsto\beta x$ (so $h_i\mapsto\beta h_i$) and $\sigma\mapsto\sigma/\beta$,
which keeps every product $\sigma h_i$ fixed; the shape of $f$ in the $(x,y)$ plane is therefore
unchanged, only horizontally stretched. Under $x\mapsto\beta x$ a derivative carries a factor
$\beta^{-1}$: $f'\mapsto\beta^{-1}f'$, $f''\mapsto\beta^{-2}f''$, while $dx\mapsto\beta\,dx$. Thus
$B=\int(f'')^2dx\mapsto\beta^{-4}\cdot\beta\,B=\beta^{-3}B$, and with $\sigma^2\mapsto\beta^{-2}\sigma^2$,
$\sigma^2 A=\sigma^2\!\int(f')^2dx\mapsto\beta^{-2}\cdot\beta^{-2}\cdot\beta\,\sigma^2A=\beta^{-3}\sigma^2A$.
Both numerator and denominator scale by $\beta^{-3}$, so $\varphi$ is unchanged. This is the
days-versus-years invariance.
\end{proof}

This removes the unit dependence that afflicts an unnormalised energy criterion, under which a
choice of basis points versus decimals, or days versus years, would silently move the optimum.
The target $\varphi^\star$ is the single risk--hedging dial called for in the introduction:
$\varphi^\star\!\to\!1$ requests a smoother, cubic-like curve; $\varphi^\star\!\to\!0$ buys
locality and forward stability at the cost of curvature. The procedure is summarised as follows.

\begin{enumerate}[leftmargin=2.2em,itemsep=1pt]
\item Choose a target fraction $\varphi^\star$ (default $\tfrac12$).
\item For trial $\sigma$, solve \eqref{eq:ttridiag} for $\{z_i\}$ and evaluate
$E_{\mathrm{bend}}(\sigma)$ and $E_{\mathrm{tens}}(\sigma)$ in closed form via
\eqref{eq:Eb} and \eqref{eq:Ea}.
\item Form $\varphi(\sigma)$ from \eqref{eq:phi}.
\item Solve $\varphi(\sigma^\star)=\varphi^\star$ by bisection (one-dimensional, monotone).
\item Return $\sigma^\star$ and the curve $f_{\sigma^\star}$.
\end{enumerate}

\begin{figure}[t]\centering
\includegraphics[width=0.74\textwidth]{fig_phi.png}
\caption{The bending-energy fraction $\varphi(\sigma)$ for the curve of Table~\ref{tab:data}.
It descends smoothly and monotonically from one (cubic) to zero (linear); equipartition
($\varphi=\tfrac12$) is attained at $\sigma^\star=4.86$. The dotted lines mark
$\varphi^\star=0.75$ ($\sigma=2.76$) and $\varphi^\star=0.25$ ($\sigma=8.69$). Contrast the
monotone single crossing here with the spiky, range-dependent energy curves produced by the
point-load construction.}\label{fig:phi}
\end{figure}

\subsection{Results}
For the Treasury curve of Table~\ref{tab:data} the criterion gives $\sigma^\star=2.76$ at
$\varphi^\star=0.75$, $\sigma^\star=4.86$ at the equipartition default, and $\sigma^\star=8.69$
at $\varphi^\star=0.25$. Figure~\ref{fig:energy} plots the two component energies. The flexural
energy is nearly flat in $\sigma$ --- the data's curvature is concentrated in the densely
knotted short end, which does not straighten until very high tension --- while the axial energy
climbs from zero and crosses $E_{\mathrm{bend}}$ exactly at the equipartition point. The
crossover is therefore driven by the growth of the tension energy against a slowly varying
bending energy, and the optimum is well separated from both spline limits.

\begin{figure}[t]\centering
\includegraphics[width=0.74\textwidth]{fig_energy.png}
\caption{Flexural energy $E_{\mathrm{bend}}$ and axial energy $E_{\mathrm{tens}}$ (log scale,
per unit $EI$) versus $\sigma$. The crossing $E_{\mathrm{bend}}=E_{\mathrm{tens}}$ is the
equipartition optimum $\sigma^\star=4.86$.}\label{fig:energy}
\end{figure}

\section{Validation}\label{sec:validation}
An interpolating spline reproduces every input instrument exactly, so the in-sample repricing
error is identically zero and cannot discriminate between values of $\sigma$. The discriminating
tests are the no-arbitrage and hedging properties that motivated the construction.

\paragraph{Forward positivity and overshoot.} A necessary no-arbitrage condition is that the
implied instantaneous forward $f(t)=y(t)+t\,y'(t)$ be non-negative. On the curve of
Table~\ref{tab:data} both the cubic and the tension interpolant keep forwards comfortably
positive (minimum $3.64\%$), as one expects of a steep upward curve; this data set does not
trigger the negative-forward pathology, and we do not manufacture one. The discriminating
feature lies elsewhere, in the long end (Figure~\ref{fig:forwards}). The flat $20$--$30$-year
par segment implies a near-flat forward of about $5.07\%$, which the tension curve reproduces.
The cubic instead fabricates a forward hump rising to almost $6\%$ near eighteen years and then
declining to $4.7\%$ at thirty years --- a shape with no support in the data, the same
overshoot that is visible directly in the yield curve of Figure~\ref{fig:curves}. The price of
the tension curve's fidelity is a kink in the forward at twenty years, where the steep
$10$--$20$-year segment meets the flat $20$--$30$-year one; this is the smoothness-versus-fidelity
trade-off that $\varphi^\star$ governs, and it is also a reminder, in the spirit of Hagan and
West~\cite{Hagan}, that when forward smoothness is itself the objective one should tension the
forward curve rather than the yield curve --- the area-constrained construction we develop in
Section~\ref{sec:fwd}.

\begin{figure}[t]\centering
\includegraphics[width=0.74\textwidth]{fig_forwards.png}
\caption{Implied instantaneous forwards. The cubic manufactures a forward hump to nearly $6\%$
and a declining long end unsupported by the flat $20$--$30$-year par data; the tension curve at
$\sigma^\star$ tracks the data-implied flat forward, at the cost of a kink at twenty
years.}\label{fig:forwards}
\end{figure}

\paragraph{Hedge locality.} The economically decisive property is locality: how far along the
curve a perturbation of a single input instrument propagates. We perturb the forward at node
$t_k$ by $\varepsilon$, rebuild the curve at fixed $\sigma$, and measure the influence
$\partial f(t)/\partial r_k$ and the \emph{locality radius}
\begin{equation}
\rho_k=\frac{\int |t-t_k|\,\big|\partial f(t)/\partial r_k\big|\,dt}
{\int \big|\partial f(t)/\partial r_k\big|\,dt},
\label{eq:rho}
\end{equation}
the mean distance of the response from the perturbed node; smaller $\rho_k$ means a more local,
more easily hedged curve. Figure~\ref{fig:locality} bumps the benchmark $5$-year node. The
cubic response stays between ten and forty percent of its peak across the \emph{entire}
thirty-year curve --- secondary lobes near thirteen and twenty-four years --- giving
$\rho_5=6.7$ years, whereas the tension spline at $\sigma^\star$ decays roughly exponentially in
$\sigma|t-t_k|$ and is effectively extinguished beyond ten years, confining the response to
$\rho_5=0.8$ year, an eight-fold improvement. The gain is consistent across the belly of the
curve: the locality ratio between cubic and tension ranges from about seven at the front to
nine at two years, narrowing only at the sparsely knotted $20$--$30$-year end, where the
ten-year gaps between nodes limit the locality attainable by any interpolant. The
energy-partition tension thus delivers precisely the local hedges that the introduction posed
as the goal, and does so through a single, interpretable parameter.

The exponential decay has a clean physical reading. The influence of a perturbation is governed
by the Green's function of the spline operator $\mathcal{L}=d^4/dx^4-\sigma^2 d^2/dx^2$, i.e.\ the
solution of $\mathcal{L}\,G(x,x_k)=\delta(x-x_k)$. Away from the source this satisfies the
homogeneous equation $G^{(\mathrm{iv})}-\sigma^2 G''=0$, whose decaying mode is $e^{-\sigma|x-x_k|}$.
Thus $\sigma$ is an inverse screening length: a rate shock at $t_k$ is an evanescent disturbance
that dies out over a distance $\sim1/\sigma$, which is exactly the hedge-locality radius
$\rho_k$. In the untensioned ($\sigma\to0$) cubic limit the operator degenerates to
$d^4/dx^4$, whose Green's function decays only algebraically and oscillates --- the long-range
``ringing'' of Figure~\ref{fig:locality} --- so a shock propagates across the whole curve like an
undamped wave. Tension plays the role of a restoring force that damps this propagation; raising
$\sigma$ shortens the screening length and localises the hedge.

\begin{figure}[t]\centering
\includegraphics[width=0.74\textwidth]{fig_locality.png}
\caption{Decay of the single-node influence $|\partial f/\partial r_k|$ (log scale) for a bump
at the benchmark $5$-year node. The cubic (locality radius $\rho=6.7$ yr) propagates the
perturbation across the whole curve; the tension spline at $\sigma^\star$ ($\rho=0.8$ yr)
localises it within a few years.}\label{fig:locality}
\end{figure}

\section{Tensioning the Instantaneous Forward Curve}\label{sec:fwd}
The validation of Section~\ref{sec:validation} exposed a structural point. Smoothness is a
property of the interpolated variable and is degraded by differentiation: if the zero curve
$r(t)$ is the $C^2$ tension spline, the instantaneous forward $f(t)=r(t)+t\,r'(t)$ is only
$C^1$, and as $\sigma$ grows and $r$ approaches piecewise-linear, $f$ develops kinks at the
knots. The three natural interpolation variables, linked by
$I(t):=-\ln P(t)=r(t)\,t=\int_0^t f$, behave differently under tension in the high-$\sigma$
limit: tensioning $r$ gives a kinked piecewise-linear forward; tensioning $I$ (the
log-discount) gives $f=I'$, a piecewise-\emph{constant} forward (the familiar local
``constant-forward'' rule); and tensioning $f$ itself keeps the forward a continuous $C^2$
spline for every $\sigma$. Hence, \emph{if forward smoothness is the objective, the forward
curve is the object to tension.}

The forward is not observed pointwise; the market pins the discount factors, equivalently the
integrals of $f$ over each interval. Tensioning the forward is therefore an \emph{area-constrained}
variational problem: minimise the forward's own tension energy subject to reproducing every
discount factor,
\begin{equation}
\min_{f}\ \int_{0}^{T}\!\Big[\tfrac12\big(f''\big)^2+\tfrac12\sigma^2\big(f'\big)^2\Big]\,dt
\quad\text{s.t.}\quad
\int_{t_{i-1}}^{t_i} f(u)\,du = I(t_i)-I(t_{i-1})=:a_i,\ \ i=1,\dots,m.
\label{eq:fwdmin}
\end{equation}
With a multiplier $\lambda_i$ for each area constraint, form the Lagrangian
\[
\mathcal{L}[f,\lambda]=\int_0^T\!\Big[\tfrac12(f'')^2+\tfrac12\sigma^2(f')^2\Big]dt
-\sum_{i=1}^m\lambda_i\Big(\int_{t_{i-1}}^{t_i}\!f\,du-a_i\Big).
\]
The constraint integral contributes $-\lambda_i$ to the variation on $(t_{i-1},t_i)$, while the
energy term varies exactly as in Proposition~\ref{prop:euler} (integrate $f''\delta f''$ by parts
twice and $\sigma^2 f'\delta f'$ once). Setting $\delta\mathcal{L}=0$ gives, on each interval, the
Euler--Lagrange equation with a piecewise-constant forcing,
\begin{equation}
f^{(\mathrm{iv})}-\sigma^2 f''=\lambda_i \quad\text{on } (t_{i-1},t_i),
\end{equation}
whose solution is the homogeneous tension spline plus the particular solution
$f_p=-\lambda_i t^2/(2\sigma^2)$ (since $f_p^{(\mathrm{iv})}=0$ and $-\sigma^2 f_p''=\lambda_i$).
The $m$ multipliers are fixed by the $m$ area constraints, together with continuity of
$f,f',f''$ across knots and the natural end conditions $f''=f'''=0$. In practice the problem is
discretised on a fine grid $\{t_k\}$: with $D_1,D_2$ the first- and second-difference operators,
the energies become the quadratic forms $\int(f')^2\approx f^{\!\top}\!D_1^{\!\top}D_1 f\,\Delta t=:f^{\!\top}\!K_t f$
and $\int(f'')^2\approx f^{\!\top}\!K_b f$, and the area constraints become $Cf=a$ with $C_{ik}$ the
length of grid cell $k$ inside interval $i$. Minimising $\tfrac12 f^{\!\top}(K_b+\sigma^2K_t)f$
subject to $Cf=a$ has the first-order (KKT) conditions $(K_b+\sigma^2K_t)f+C^{\!\top}\lambda=0$ and
$Cf=a$, i.e.\ the single sparse symmetric system
\begin{equation}
\begin{pmatrix} K_b+\sigma^2 K_t & C^{\!\top}\\ C & 0\end{pmatrix}
\begin{pmatrix} f\\ \lambda\end{pmatrix}=\begin{pmatrix}0\\ a\end{pmatrix}.
\end{equation}
The energy-partition criterion carries over verbatim, with $\varphi(\sigma)$ now formed from
the forward's flexural energy $\int(f'')^2$ and axial energy $\sigma^2\int(f')^2$; tensioning
the forward thus balances forward \emph{curvature} against forward \emph{slope} directly.

On the curve of Table~\ref{tab:data} this construction reproduces every discount factor to
$2\times10^{-10}$ (machine precision; no-arbitrage is exact), keeps the forward positive, and
returns a forward-space equipartition tension $\sigma^\star=6.78$. Figure~\ref{fig:fwdspace}
compares, at a common tension, the forward implied by tensioning the yield curve with the
forward obtained by tensioning the forward directly. The two share the front-end wiggle, which
is genuine signal forced by the kinked short-end par data, but in the long end, where the data
are smooth, the difference is stark: the yield-tensioned forward cliffs at twenty years, whereas
the forward-tensioned curve is smooth there. Quantitatively, over $t>8$ years the
forward-tensioned curve has $1{,}285$ times lower roughness $\int(f'')^2$ and a peak curvature
at twenty years smaller by a factor of $338$, at no cost in repricing. When smooth forwards
matter --- for forward-rate-dependent payoffs, convexity adjustments, or volatility-surface
inputs --- this is the construction to use; the energy-partition criterion selects its tension
exactly as before.

\begin{figure}[t]\centering
\includegraphics[width=0.78\textwidth]{fig_fwdspace.png}
\caption{Instantaneous forward at a common tension $\sigma=4.9$. Tensioning the yield curve
(green) produces a forward that cliffs at twenty years; tensioning the forward directly (red)
keeps it smooth while reproducing all discount factors exactly. The cubic-on-yield forward
(dotted) overshoots. All three agree on the data-driven front-end wiggle.}\label{fig:fwdspace}
\end{figure}

\section{Non-Uniform Tension Calibration}\label{sec:nonunif}
A single global tension treats every segment alike, yet the segments are not: on the curve of
Table~\ref{tab:data} the interval widths $h_i$ span more than two orders of magnitude, from one
month to ten years. The dimensionless group that fixes a segment's shape is the product
$\sigma h_i$, the argument of the hyperbolic functions in \eqref{eq:tspline}; under a uniform
$\sigma$ this product ranges from $0.8$ on the shortest segment to $48.6$ on the
$20$--$30$-year segment. The short end is therefore barely tensioned while the long end is
pulled essentially straight.

The imbalance shows cleanly in the per-segment energy fraction
$\varphi_i = E^{(i)}_{\mathrm{bend}}/(E^{(i)}_{\mathrm{bend}}+E^{(i)}_{\mathrm{tens}})$, the
local analogue of \eqref{eq:phi}. At the global optimum $\sigma^\star=4.86$ the fractions run
from $0.86$ on the shortest segment down to below $0.01$ on the segments between three and
twenty years (Figure~\ref{fig:nuphi}, green): the belly and long end carry almost no flexural
energy --- they have been flattened to near-linearity --- while the short end stays curved.
Equipartition holds only on average, not segment by segment.

The remedy is to let the tension vary by segment. We assign a tension $\sigma_i$ to each
interval (the non-uniform spline of Appendix~\ref{app:nonunif}) and calibrate the vector
$\{\sigma_i\}$ so that every segment individually attains the target fraction,
\begin{equation}
\varphi_i(\sigma_i)=\varphi^\star\qquad\text{for all } i.
\label{eq:segphi}
\end{equation}
Because the per-segment energies \eqref{eq:Eb}--\eqref{eq:Ea} and the knot second derivatives
$\{z_i\}$ are coupled through the tridiagonal system, \eqref{eq:segphi} is solved by a short
fixed-point iteration: holding $\{z_i\}$ fixed, each one-dimensional equation
$\varphi_i(\sigma_i)=\varphi^\star$ is solved by bisection (each $\varphi_i$ is monotone in its
$\sigma_i$), the non-uniform tridiagonal is then re-solved for $\{z_i\}$, and the step is
repeated. On the data of Table~\ref{tab:data} the iteration converges in about forty steps and
drives every $\varphi_i$ to $\tfrac12$ to machine precision (Figure~\ref{fig:nuphi}, red).

\begin{figure}[t]\centering
\includegraphics[width=0.74\textwidth]{fig_nu_phi.png}
\caption{Per-segment energy fraction $\varphi_i$. A uniform tension (green) leaves the
fractions scattered from $0.86$ to below $0.01$; per-segment calibration (red) equalises them
to $\tfrac12$ on every segment.}\label{fig:nuphi}
\end{figure}

The calibrated tensions (Figure~\ref{fig:nusigma}) scale roughly as $\sigma_i\sim 1/h_i$, from
$\sigma\approx 12$ on the shortest segment to $\sigma\approx 0.14$ in the belly, so that the
product $\sigma_i h_i$ is compressed from the uniform $60\!:\!1$ spread into a narrow band of
order unity across the active part of the curve. Calibration thus discovers that the natural
tension scale is local: each segment is tensioned in proportion to its own width and curvature
rather than by a single number imposed on the whole curve. The resulting curve
(Figure~\ref{fig:nucurve}) restores the gentle curvature to the five-to-twenty-year belly that
the uniform spline had straightened, while still suppressing the cubic's long-end hump. It
remains $C^2$ and reprices every node exactly, since calibration changes only the tension
profile and not the interpolation conditions.

\begin{figure}[t]\centering
\includegraphics[width=0.74\textwidth]{fig_nu_sigma.png}
\caption{Calibrated per-segment tension $\sigma_i$ (red, log scale) against the uniform optimum
(green dashed), and the dimensionless product $\sigma_i h_i$ (blue). Calibration raises the
short-end tension and lowers the long-end tension, bringing $\sigma_i h_i$ into a narrow
band.}\label{fig:nusigma}
\end{figure}

\begin{figure}[t]\centering
\includegraphics[width=0.74\textwidth]{fig_nu_curve.png}
\caption{Uniform versus per-segment-calibrated tension spline. The calibrated curve recovers
curvature in the belly that the uniform spline flattens, while both avoid the cubic's long-end
hump.}\label{fig:nucurve}
\end{figure}

\section{Synthesis: Non-Uniform Tension on the Forward Curve}\label{sec:fwdnu}
The two extensions just developed are orthogonal and combine. Section~\ref{sec:fwd} makes the
instantaneous forward the tensioned object, under area constraints that reproduce the discount
factors; Section~\ref{sec:nonunif} lets the tension vary by segment to balance energy locally.
Applying the second to the first gives a \emph{non-uniform tension on the forward curve}: the
forward solves the area-constrained problem \eqref{eq:fwdmin} with a piecewise tension
$\sigma(t)=\sigma_i$ on each interval, and the $\{\sigma_i\}$ are calibrated to per-segment
equipartition of the forward's own energy.

The forward formulation makes the calibration markedly simpler than its yield-curve counterpart.
In the discretised problem the forward $f$ is the direct unknown, so for a \emph{fixed} $f$ the
per-segment flexural and (unweighted) axial energies
$B_i=\tfrac12\!\int_{t_{i-1}}^{t_i}(f'')^2$ and $A_i=\tfrac12\!\int_{t_{i-1}}^{t_i}(f')^2$ are
quadratic forms in $f$ that do not involve $\sigma$ at all; the tension enters the segment's
axial energy only through the multiplicative factor $\sigma_i^2$, as
$E^{(i)}_{\mathrm{tens}}=\sigma_i^2 A_i$, while $E^{(i)}_{\mathrm{bend}}=B_i$. Write the
per-segment fraction explicitly,
\[
\varphi_i(\sigma_i)=\frac{E^{(i)}_{\mathrm{bend}}}{E^{(i)}_{\mathrm{bend}}+E^{(i)}_{\mathrm{tens}}}
=\frac{B_i}{B_i+\sigma_i^2 A_i},
\]
and set it to the target $\tfrac12$. Then $B_i+\sigma_i^2 A_i=2B_i$, i.e.\ $\sigma_i^2 A_i=B_i$,
so
\begin{equation}
\sigma_i^{\,2}A_i=B_i\qquad\Longrightarrow\qquad \sigma_i=\sqrt{B_i/A_i},
\label{eq:sigclosed}
\end{equation}
a closed-form update with no bisection required (contrast the yield-curve case of
Section~\ref{sec:nonunif}, where the $z_i$ depend on $\sigma_i$ through the hyperbolic energies
\eqref{eq:Eb}--\eqref{eq:Ea} and a one-dimensional root-find is needed per segment). The
construction is a short fixed point: solve the KKT system for
$f$ at the current $\{\sigma_i\}$, recompute $\{B_i,A_i\}$, update each $\sigma_i$ by
\eqref{eq:sigclosed}, and repeat. On the curve of Table~\ref{tab:data} it converges in about
forty iterations, reproduces every discount factor to $4\times10^{-11}$, keeps the forward
positive, and drives all eleven per-segment fractions to $\tfrac12$, from the scatter of
$0.001$ to $0.88$ left by a uniform forward tension (Figure~\ref{fig:fwdnu}, right). The
calibrated tensions are large at the curved front end ($\sigma\approx 20$ on the
$0.08$--$0.25$-year segment) and small on the flat long end ($\sigma\approx 0.09$ on
$20$--$30$ years).

\begin{figure}[t]\centering
\includegraphics[width=0.96\textwidth]{fig_fwd_nonunif.png}
\caption{Non-uniform tension on the forward curve. \emph{Left:} the uniform and
per-segment-calibrated forwards, both smooth and both reproducing all discount factors exactly.
\emph{Right:} per-segment forward energy fractions; calibration replaces the uniform tension's
scatter with exact equipartition on every segment.}\label{fig:fwdnu}
\end{figure}

One feature of this combination deserves an honest statement. Per-segment equipartition is an
energy-\emph{balance} rule, not a forward-\emph{smoothness} rule. On a near-flat long-end
segment both $B_i$ and $A_i$ are small, and \eqref{eq:sigclosed} returns a \emph{low} tension,
which lets the forward curve freely over that interval rather than holding it taut
(Figure~\ref{fig:fwdnu}, left, beyond twenty years). This is the correct behaviour for energy
balance, but a desk that prizes a stable long-end forward would impose a lower $\varphi^\star$
there --- more tension, a straighter forward --- exactly the risk--preference dial of
Section~\ref{sec:energy}, now exercised per segment. The unified construction thus delivers a
$C^2$, exactly-repricing, energy-balanced forward curve in which a single interpretable
parameter, applied globally or segment by segment, governs the whole trade-off between
smoothness, locality, and long-end stability.

\section{Multi-Curve Construction: SOFR Discounting and a Tensioned Projection Curve}\label{sec:multicurve}
Modern fixed-income desks operate in a multi-curve world: cashflows are discounted on an
overnight-indexed (OIS / SOFR) curve, while forward rates are read off one or more separate
projection curves that may carry a tenor or credit basis \cite{Bianchetti,Henrard,Ametrano}. The interpolation method of this paper
is deliberately agnostic to that architecture --- it acts \emph{downstream} of the bootstrap ---
and we now show it in place, with OIS discounting held exact and the energy-partition tension
applied only to the projection curve.

\paragraph{Bootstrapping the OIS discount curve.} Let $S_n$ be the par rate of the OIS swap
maturing at $T_n=n$ years (annual, unit accrual). The floating leg of an OIS pays compounded
overnight rate and has present value $1-P(T_n)$, while the fixed leg has value
$S_n\sum_{i=1}^n P(T_i)$. Par means these are equal:
\begin{equation}
S_n\sum_{i=1}^{n}P(T_i)=1-P(T_n).
\label{eq:par}
\end{equation}
Writing the annuity $A_{n-1}=\sum_{i=1}^{n-1}P(T_i)$ and isolating the unknown $P(T_n)$,
\[
S_n\big(A_{n-1}+P(T_n)\big)=1-P(T_n)
\;\Longrightarrow\;
P(T_n)\big(1+S_n\big)=1-S_nA_{n-1}
\;\Longrightarrow\;
P(T_n)=\frac{1-S_nA_{n-1}}{1+S_n},
\]
a forward recursion seeded by $A_0=0$, i.e.\ $P(T_1)=1/(1+S_1)$ \cite{Ametrano}. This produces discount factors
that reprice every input swap exactly; on the representative SOFR levels used here the maximum
reprice error is $6\times10^{-17}$. This bootstrapped curve is the \emph{crisp} discounting
curve and is left untouched.

\paragraph{Tensioning the projection curve.} The projection curve carries a (small, illustrative)
tenor basis $b(T)$ over the OIS zeros, $z_{\mathrm{proj}}(T_i)=z_{\mathrm{OIS}}(T_i)+b(T_i)$ with
$z=-\ln P/T$. We build its instantaneous forward by the area-constrained construction of
Section~\ref{sec:fwd}: the forward is a tension spline whose integral over each interval
reproduces the projection discount factors, and $\sigma$ is fixed by the energy-partition
criterion $\varphi(\sigma^\star)=\tfrac12$. On the data here this gives $\sigma^\star=0.73$, a
projection forward that reproduces the projection discount factors to $1.5\times10^{-10}$ and
stays strictly positive (minimum $3.62\%$), and that is visibly smoother than the
cubic-in-yield forward, which overshoots at the long end where the par curve is humped by the
negative $20$--$30$-year swap spread (Figure~\ref{fig:multicurve}).

\begin{figure}[t]\centering
\includegraphics[width=0.98\textwidth]{fig_multicurve.png}
\caption{Multi-curve SOFR/OIS construction. \emph{Left:} the bootstrapped OIS discount zero curve
(crisp, exact at every node) and the projection-curve nodes offset by the tenor basis.
\emph{Right:} the projection instantaneous forward built by the area-constrained energy-partition
tension spline (smooth, positive, reproducing the projection discount factors exactly), the
cubic-in-yield forward (overshooting at the long end), and the OIS forward for reference. SOFR
levels are representative of the 22-May-2026 environment; the basis is illustrative; the
discounting bootstrap and discount-factor reproduction are exact.}\label{fig:multicurve}
\end{figure}

The example makes the architectural point concrete. Nothing in the energy-partition machinery
touches the discounting bootstrap: the OIS curve is fixed by \eqref{eq:par}, and the tension
spline, its energy criterion, and (if desired) the per-segment calibration of
Section~\ref{sec:nonunif} all live entirely on the projection leg. The same construction would
apply unchanged to a cross-currency or credit-sensitive projection curve. We are transparent
that, because live mid-market SOFR swap quotes are subscription-gated, the par levels here are
representative of the period rather than a specific fixing, and the tenor basis is a stylised
ramp; the bootstrap recursion and the discount-factor reproduction are exact regardless of the
levels chosen.

\section{Risk Sensitivities and Stability of the Dial}\label{sec:risk}
A construction is only usable on a desk if the sensitivities of the curve to its inputs are
available cheaply and if the selected tension does not lurch as the market moves. Both follow
analytically from the structure already in place.

\paragraph{The curve delta is analytic.} The knot second derivatives solve the tridiagonal
system $M(\sigma)\,z=b(y)$ of \eqref{eq:ttridiag}, in which the matrix $M$ depends only on
$\sigma$ and the mesh while the right-hand side $b$ is \emph{linear} in the input rates,
$b_i=\gamma_i-\gamma_{i-1}$ with $\gamma_k=\sigma^2(y_{k+1}-y_k)/h_k$. Differentiating,
\begin{equation}
\frac{\partial z}{\partial y_j}=M^{-1}\frac{\partial b}{\partial y_j},
\qquad
\frac{\partial b_i}{\partial y_j}=\frac{\sigma^2}{h_i}(\delta_{j,i+1}-\delta_{j,i})
-\frac{\sigma^2}{h_{i-1}}(\delta_{j,i}-\delta_{j,i-1}),
\label{eq:dzdy}
\end{equation}
which costs one back-substitution per node against the already-factored $M$, i.e.\ $O(N)$ for the
whole Jacobian. Since the interpolant \eqref{eq:tspline} is linear in $(y,z)$, the curve delta
$\partial f(x)/\partial y_j$ is then closed form: differentiating \eqref{eq:tspline} term by term,
\begin{equation}
\frac{\partial f}{\partial y_j}
=\frac{\partial_j z_i\,\sinh(\sigma(x_{i+1}-x))+\partial_j z_{i+1}\,\sinh(\sigma(x-x_i))}{\sigma^2\sinh(\sigma h_i)}
+\Big(\delta_{ji}-\tfrac{\partial_j z_i}{\sigma^2}\Big)\frac{x_{i+1}-x}{h_i}
+\Big(\delta_{j,i+1}-\tfrac{\partial_j z_{i+1}}{\sigma^2}\Big)\frac{x-x_i}{h_i},
\end{equation}
where $\partial_j z\equiv\partial z/\partial y_j$. This is the exact hedge Jacobian; the
finite-difference bump experiment of Section~\ref{sec:validation} is precisely its numerical
shadow, and the localised triangular response seen there is the column $\partial f/\partial y_j$.
The analytic and finite-difference Jacobians agree to $2\times10^{-11}$
(Figure~\ref{fig:jacobian}, left). A desk thus obtains delta and gamma against every node without
re-bootstrapping or bumping.

\paragraph{The dial is a smooth function of the inputs.} The selected tension solves
$\varphi(\sigma^\star;y)=\varphi^\star$. By the implicit function theorem,
\begin{equation}
\frac{\partial\sigma^\star}{\partial y_j}
=-\frac{\partial\varphi/\partial y_j}{\partial\varphi/\partial\sigma}\Bigg|_{\sigma=\sigma^\star},
\label{eq:dsigstar}
\end{equation}
where $\partial\varphi/\partial y_j$ follows from the closed-form energies \eqref{eq:Eb}--\eqref{eq:Ea}
and \eqref{eq:dzdy}, and $\partial\varphi/\partial\sigma$ is the slope established in
Proposition~\ref{prop:phi}. The right-hand side is finite and continuous \emph{exactly where}
$\partial\varphi/\partial\sigma\neq0$ --- the same strict-monotonicity condition that makes
$\sigma^\star$ unique. Hence $\sigma^\star$ is a $C^1$ function of the input rates: it cannot jump.
On the curve of Table~\ref{tab:data}, $\partial\varphi/\partial\sigma=-0.099\neq0$, and
\eqref{eq:dsigstar} reproduces the directly re-optimised $\sigma^\star$ to a relative error of
$6\times10^{-5}$.

\begin{figure}[t]\centering
\includegraphics[width=0.98\textwidth]{fig_jacobian.png}
\caption{Risk sensitivities and stability. \emph{Left:} the analytic curve delta
$\partial f/\partial y_j$ (lines) for the five- and ten-year nodes, overlaid with finite
differences (circles); the response is local and the two agree to $2\times10^{-11}$.
\emph{Right:} the equipartition $\sigma^\star$ under $\pm100$\,bp parametric stress. It is exactly
invariant to a parallel level shift (which moves the curve but not its derivatives) and varies
smoothly and monotonically under slope and curvature shocks --- no regime jumps.}\label{fig:jacobian}
\end{figure}

Figure~\ref{fig:jacobian} (right) confirms the stability empirically. A parallel level shift
leaves $\sigma^\star$ \emph{exactly} unchanged, because adding a constant to every rate shifts the
curve vertically without altering $f'$ or $f''$, hence without altering either energy; slope and
curvature shocks of $\pm100$ basis points move $\sigma^\star$ smoothly and monotonically across a
bounded range. The dial is therefore safe to hold fixed intraday and to update slowly.

Finally, a remark on the multi-curve setting of Section~\ref{sec:multicurve}: because the
discount curve is bootstrapped independently and held fixed, the projection-curve sensitivities
do not feed back into it. The cross-curve Jacobian is block-diagonal --- tensioning the
projection leg leaves the discount factors untouched by construction --- so no additional
cross-curve risk term arises from the energy-partition step.

\section{Empirical Comparison and the Role of $\varphi^\star$}\label{sec:empirical}
This section situates the method quantitatively against the interpolators a desk would actually
consider, characterises the effect of the risk dial $\varphi^\star$, and probes robustness under
stressed curve shapes. All numbers are computed on the real par-yield curve of
Table~\ref{tab:data}; the benchmark methods are implemented from their published definitions.
We are deliberately even-handed: the tension spline is not best on every axis, and the table
makes clear precisely what it does and does not buy.

\paragraph{Benchmarks.} We compare against five alternatives: the natural cubic; the Hagan--West
monotone-convex forward scheme \cite{Hagan}; Smith--Wilson with an illustrative ultimate forward
rate ($3.9\%$) and mean-reversion $\alpha=0.10$; log-linear discount-factor interpolation
(equivalently a piecewise-constant instantaneous forward); and an explicit piecewise-flat forward.
For each we measure the forward roughness $\int (f_{\mathrm{fwd}}')^2\,dt$ (a global smoothness
penalty on the instantaneous forward), the minimum implied forward (positivity), and, where the
method is a local interpolant for which a node bump is well defined, the $5$-year hedge-locality
radius of Section~\ref{sec:validation}.

\begin{table}[h]\centering\small
\begin{tabular}{lccc}
\hline
Method & Forward roughness ($\times10^{4}$) & Min.\ implied forward & $5$y locality (yr)\\
\hline
Natural cubic                       & $2.03$  & $3.64\%$ & $6.66$\\
Tension, $\varphi^\star=\tfrac12$   & $2.84$  & $3.64\%$ & $\mathbf{0.79}$\\
Hagan--West monotone convex         & $1.84$  & $3.63\%$ & ---\\
Smith--Wilson (UFR/$\alpha$ illustr.)& $1.44$ & $3.63\%$ & ---\\
Log-discount (log-linear DF)        & $23.94$ & $3.66\%$ & ---\\
Piecewise-flat forward              & $23.94$ & $3.66\%$ & ---\\
\hline
\end{tabular}
\caption{Benchmark interpolation methods on the $22$-May-$2026$ curve. Forward roughness is the
integrated squared slope of the instantaneous forward; smaller is smoother. All methods keep the
implied forward positive on this upward curve, so positivity is not a discriminator here. The
tension spline's decisive advantage is locality, not global smoothness.}\label{tab:bench}
\end{table}

Three honest conclusions follow. First, \emph{positivity is not a discriminator on this curve}:
every method, including the cubic, keeps the implied forward comfortably positive, because the
underlying data are upward-sloping. The well-known negative-forward pathology of the cubic
appears on differently shaped curves, not this one, and we do not manufacture such a curve.
Second, the tension spline is \emph{not} the globally smoothest forward: Smith--Wilson and
Hagan--West both achieve lower forward roughness, because they are designed to optimise forward
shape directly, while the piecewise-flat and log-discount methods are an order of magnitude
rougher (their forwards are kinked or stepped). Third, and decisively, the local interpolants
separate sharply on \emph{locality}: a $5$-year node bump propagates over $6.66$ years under the
cubic but only $0.79$ year under the equipartition tension spline --- an eight-fold compression
that the globally smooth methods do not provide, because Smith--Wilson and monotone-convex
schemes are not local interpolators in the node-bump sense. The tension spline thus occupies a
distinct point in the design space: a local, $C^2$ interpolant whose locality is tunable, rather
than a global shape-optimiser.

\paragraph{The risk dial $\varphi^\star$.} The default $\varphi^\star=\tfrac12$ is a neutral
choice, not a claimed optimum: there is no curve-construction optimum independent of a desk's
loss function, and the whole point of $\varphi$ is to expose the trade-off rather than hide it.
Table~\ref{tab:phisweep} characterises the dial. Lowering $\varphi^\star$ (more tension)
tightens locality and smooths the forward at the cost of pulling the curve toward the linear
shape; raising it does the reverse. The dependence is smooth and monotone --- the desk reads off
the setting that matches its own smoothness-versus-locality preference.

\begin{table}[h]\centering\small
\begin{tabular}{ccccc}
\hline
$\varphi^\star$ & $\sigma^\star$ & Forward roughness ($\times10^4$) & $5$y locality (yr) & Min.\ fwd\\
\hline
$0.3$ & $7.60$ & $3.53$ & $0.74$ & $3.64\%$\\
$0.4$ & $6.01$ & $3.12$ & $0.76$ & $3.64\%$\\
$0.5$ & $4.86$ & $2.84$ & $0.79$ & $3.64\%$\\
$0.6$ & $3.94$ & $2.63$ & $0.84$ & $3.64\%$\\
$0.7$ & $3.14$ & $2.46$ & $0.91$ & $3.64\%$\\
\hline
\end{tabular}
\caption{Characterisation of the risk dial on the $22$-May-$2026$ curve. As $\varphi^\star$ rises,
$\sigma^\star$ falls (less tension), the forward becomes smoother, and the hedge-locality radius
widens. The monotone, gradual dependence is what makes $\varphi^\star$ a usable control rather
than a sensitive knob.}\label{tab:phisweep}
\end{table}

\begin{figure}[t]\centering
\includegraphics[width=0.98\textwidth]{fig_benchmark.png}
\caption{\emph{Left:} instantaneous forwards from all six methods on the $22$-May-$2026$ curve;
the piecewise-flat and log-discount forwards are stepped, the cubic and tension splines and the
two global schemes are smooth. \emph{Right:} the tension forward as the dial $\varphi^\star$
sweeps from $0.3$ to $0.7$; the curves are visually close because this curve's shape is dominated
by the long-end feature, while the locality differences of Table~\ref{tab:phisweep} are the
operative distinction.}\label{fig:bench}
\end{figure}

\paragraph{Stress shapes.} Finally we re-fit the equipartition criterion to stressed transforms
of the real curve: an inverted curve (the same levels in reverse order), a steepened curve
($+150$ basis points ramped onto the long end), and a sparse curve (the $7$- and $20$-year nodes
removed). In every case the solver returns a well-defined $\sigma^\star$ and a strictly positive
implied forward --- $\sigma^\star=6.35$ (inverted, min.\ forward $2.98\%$), $3.39$ (steep,
$3.66\%$), and $5.02$ (sparse, $3.64\%$) --- confirming that the construction degrades gracefully
away from the base case.

\paragraph{Historical study, 2006--2026.} To move beyond a single date we fit the equipartition
criterion to a long historical sample: $1{,}064$ complete weekly U.S. Treasury par-yield curves
from 8 January 2006 to 24 May 2026, retrieved from the Federal Reserve (FRED) constant-maturity
series, spanning the global financial crisis, the zero-interest-rate (ZIRP) period, the
2015--2018 and 2022--2023 hiking cycles, the COVID shock, and the 2023--2024 inversion. Every
curve is a real fixing; nothing is simulated. The accompanying retrieval module reproduces the
sample, and the analysis below, from a single command. Three findings stand out
(Figure~\ref{fig:history}, Table~\ref{tab:history}).

First, the selected tension \emph{tracks the rate regime but does not jump}. The median
$\sigma^\star$ falls from about $4$ in the steep pre-crisis curves of 2006--2007 to about $0.5$
during ZIRP --- flat curves need little tension to balance their energy --- and rises again
through the subsequent hiking cycles, with the full sample ranging from $0.2$ to $8.5$. Yet the
week-over-week change is small: the median $|\Delta\sigma^\star|$ is $0.34$, with a $95$th
percentile of $2.2$. This is the empirical counterpart of the $C^1$ stability proved in
Section~\ref{sec:risk}: $\sigma^\star$ moves continuously with the market rather than lurching,
across twenty years and every regime in the sample.

Second, and most strikingly, the locality--tension law $\rho_5\sim1/\sigma^\star$ predicted by the
Green's-function argument of Section~\ref{sec:validation} holds across the \emph{entire} history.
The correlation between $\log\sigma^\star$ and $\log\rho_5$ over all $1{,}064$ weeks is $-0.987$,
and the weekly points collapse onto a single curve $\rho_5\approx 2.18/\sigma^\star$ regardless of
year or regime (Figure~\ref{fig:history}b). The inverse-screening-length interpretation, derived
on one curve, is thus confirmed out of sample on two decades of real data.

Third, the construction \emph{almost always} produces an arbitrage-free forward without any
positivity constraint being imposed: only $2.4\%$ of the $1{,}064$ weeks exhibit any negative
implied forward, and the most negative value over the whole history is $-0.008\%$ --- a
rounding-level dip confined to the post-crisis and ZIRP years when the front end was pinned near
zero. Separately, under a single global $\sigma^\star$ the per-segment energy fraction spans a
median range of $0.92$ (from near $0$ to near $0.9$) across the sample, confirming on real data
the segment-level imbalance that motivates the per-segment calibration of
Section~\ref{sec:nonunif}.

\begin{figure}[t]\centering
\includegraphics[width=0.98\textwidth]{fig_history.png}
\caption{Historical study over $1{,}064$ weekly U.S. Treasury curves, 2006--2026 (FRED).
\emph{(a)} The equipartition tension $\sigma^\star$ through time, with rate regimes shaded; it
tracks the regime smoothly. \emph{(b)} The $5$-year hedge-locality radius against $\sigma^\star$;
all weeks, coloured by year, collapse onto $\rho_5\approx c/\sigma^\star$ (log-correlation
$-0.99$), confirming the inverse-screening-length law of Section~\ref{sec:validation} out of
sample. \emph{(c)} The distribution of $\sigma^\star$ by year.}\label{fig:history}
\end{figure}

\begin{table}[h]\centering\small
\begin{tabular}{lccccc}
\hline
Regime & Period & Median $\sigma^\star$ & Median $\rho_5$ (yr) & Inverted & Neg.\ fwd.\\
\hline
Pre-crisis      & 2006--2007 & $3.95$ & $0.84$ & $30\%$ & $0\%$\\
Crisis          & 2008       & $2.37$ & $1.05$ & $0\%$  & $8\%$\\
ZIRP            & 2009--2014 & $0.53$ & $3.67$ & $0\%$  & $3\%$\\
Hiking          & 2015--2018 & $1.43$ & $1.54$ & $0\%$  & $5\%$\\
COVID / recovery& 2020--2021 & $0.77$ & $2.71$ & $0\%$  & $2\%$\\
Hiking / inversion & 2022--2024 & $1.93$ & $1.20$ & $69\%$ & $0\%$\\
\hline
Full sample     & 2006--2026 & $1.35$ & $1.61$ & $13\%$ & $2.4\%$\\
\hline
\end{tabular}
\caption{Equipartition tension and hedge locality by rate regime, $1{,}064$ weekly U.S. Treasury
curves. ``Inverted'' is the share of weeks with the $1$-month yield above the $30$-year (a coarse
proxy); ``Neg.\ fwd.'' is the share with any negative implied forward. The tension falls in flat
regimes and rises in steep ones, while the locality radius moves inversely, exactly as the theory
predicts.}\label{tab:history}
\end{table}

\paragraph{Benchmarks across regimes.} The single-curve comparison of Table~\ref{tab:bench}
raises the obvious question of whether its conclusions survive across market environments. We
therefore re-ran all six interpolators on each of the $1{,}064$ historical curves and recorded
the distribution of each metric (Table~\ref{tab:benchhist}, Figure~\ref{fig:benchhist}). The
single-curve picture holds in distribution. On \emph{forward smoothness} the tension spline sits
in the same tight band as the cubic, Hagan--West, and Smith--Wilson (median roughness
$1.9\times10^{-4}$, against $1.8$, $1.9$ and $1.6$ respectively), while the log-discount and
piecewise-flat forwards are an order of magnitude rougher in every percentile --- so tension is
competitive on smoothness without being the smoothest, exactly as on the single curve. On
\emph{locality}, the separation is decisive and universal: the tension spline's $5$-year
locality radius is strictly tighter than the cubic's in \emph{every one} of the $1{,}064$ weeks,
by a median factor of $4.1$ (and from $1.4\times$ to $8.5\times$ across the $5$th--$95$th
percentiles). The cubic's radius is a flat $6.66$ years on every curve --- its node-bump response
is scale-invariant --- whereas the tension radius ranges from $0.8$ to $4.9$ years and is almost
always far smaller (Figure~\ref{fig:benchhist}b). On \emph{positivity}, all the local
interpolants behave similarly (each shows a negative implied forward in $2$--$4\%$ of weeks, all
rounding-level and confined to the near-zero-rate years), while the stepped forwards are positive
more often only because they are piecewise constant.

\begin{table}[h]\centering\small
\begin{tabular}{lccccc}
\hline
& \multicolumn{3}{c}{Forward roughness ($\times10^4$)} & \multicolumn{2}{c}{$5$y locality (yr)}\\
Method & p5 & median & p95 & median & neg.\ fwd.\\
\hline
Natural cubic                     & $0.38$ & $1.77$ & $6.85$ & $6.66$ & $2.4\%$\\
Tension ($\varphi^\star=\tfrac12$)& $0.41$ & $1.88$ & $7.70$ & $\mathbf{1.61}$ & $2.4\%$\\
Hagan--West                       & $0.39$ & $1.92$ & $7.19$ & --- & $4.0\%$\\
Smith--Wilson                     & $0.41$ & $1.64$ & $6.54$ & --- & $3.7\%$\\
Log-discount                      & $4.16$ & $23.89$& $73.72$& --- & $0.5\%$\\
Piecewise-flat forward            & $4.16$ & $23.89$& $73.72$& --- & $0.5\%$\\
\hline
\end{tabular}
\caption{Distribution of each metric over $1{,}064$ weekly curves, 2006--2026. Forward-roughness
percentiles show the tension spline in the same band as the cubic and the two global schemes, and
far below the stepped methods. The cubic's locality radius is constant at $6.66$ years (its
node-bump response is scale-invariant); the tension radius has median $1.61$ years and is strictly
smaller than the cubic's in $100\%$ of weeks.}\label{tab:benchhist}
\end{table}

\begin{figure}[t]\centering
\includegraphics[width=0.98\textwidth]{fig_bench_history.png}
\caption{Cross-regime benchmark distributions over $1{,}064$ weekly curves. \emph{(a)} Forward
roughness by method (log scale): the tension spline (red) lies with the cubic, Hagan--West and
Smith--Wilson, an order of magnitude below the log-discount and piecewise-flat forwards.
\emph{(b)} The $5$-year hedge-locality radius: the tension distribution lies entirely to the left
of the cubic's constant $6.66$-year value, so tension is more local on every curve in the
sample.}\label{fig:benchhist}
\end{figure}

\section{Why an Energy-Partition Tension Spline? Benefits and Applications}\label{sec:why}
A fair question confronts any new curve-construction proposal: the market already has good
interpolators, so why this one? We answer it directly, situating the method against current
practice and being explicit about both its advantages and the cases in which a practitioner
should reach for something else.

\subsection{The competitive landscape}
Four families dominate practice. \emph{Simple interpolation} --- linear on rates or log-linear
on discount factors --- is robust and local but produces kinked, non-differentiable forwards.
The \emph{monotone convex} method of Hagan and West~\cite{Hagan} guarantees positive,
shape-preserving forwards and is the method the U.S. Treasury itself adopted in 2021
\cite{TreasuryMethod}; its forwards are, however, only piecewise-defined and can exhibit flat
spots, and it offers no continuous control over smoothness. The \emph{Smith--Wilson} method,
mandated for Solvency~II discounting, is in fact a close cousin of the present approach: it too
is the solution of a tensioned variational problem, and its decay parameter $\alpha$ plays
exactly the role of our $\sigma$ --- a point that lends the energy interpretation immediate
regulatory familiarity. Finally, parametric \emph{Nelson--Siegel--Svensson} models, used by
central banks to publish parsimonious curves, fit four to six economically interpretable
factors but do not reprice every instrument and are smoothing, not interpolating, devices.

The energy-partition tension spline is not a wholesale replacement for these. It is an exact
interpolant, like simple and monotone-convex interpolation and unlike NSS, but it adds what
those interpolants lack: a single continuous parameter, with a transparent physical meaning,
that moves the curve along the entire continuum from the smooth cubic to the local,
shape-preserving linear limit, together with a principled, automatic rule for placing it.

\subsection{What is distinctive}
\paragraph{One model, one auditable parameter.} Rather than choosing \emph{between} methods, a
desk chooses a \emph{point} on a one-parameter family that contains the cubic and the linear
interpolant as its endpoints. For model governance this is decisive: a single model with a
single parameter $\varphi^\star$ is far easier to document, validate, and reproduce than a
patchwork of interpolation rules, and the parameter's meaning --- the fraction of stored
elastic energy carried by curvature --- is explainable to a risk committee in one sentence.

\paragraph{A principled, scale-invariant selection rule.} The chief obstacle to adopting
tension splines has always been the absence of a defensible way to choose the tension; existing
automatic rules~\cite{Renka} optimise constraints that do not map to a desk's preferences, and
ad hoc energy rules depend on the units and the search range. The bending-energy fraction
\eqref{eq:phi} is intrinsic, dimensionless, scale-invariant and monotone, so it returns a unique
$\sigma^\star$ that does not move when rates are quoted in basis points rather than decimals.
This selection rule, not the tension spline itself, is the paper's contribution.

\paragraph{$C^2$ smoothness with controlled overshoot.} The tension spline is twice
continuously differentiable everywhere, so the curve and --- when one tensions the forward
curve --- the forwards are smooth, unlike the merely piecewise forwards of monotone-convex or
linear schemes. Yet, as Figures~\ref{fig:curves} and~\ref{fig:forwards} show, raising the
tension suppresses the spurious humps a cubic invents from flat data. The method thus occupies
a genuine middle ground: smoothness when wanted, fidelity and locality when needed, selected by
one dial rather than by switching models.

\paragraph{Risk preference as policy, not accident.} Because $\varphi^\star$ is explicit, the
trade-off between a smooth curve and local hedges is set deliberately, as desk or firm policy,
and applied consistently across curves and dates, rather than emerging as an uncontrolled
by-product of an interpolation choice.

\subsection{Applications}
The locality result of Section~\ref{sec:validation} is the most directly monetisable benefit.
Because a perturbation of one input decays within roughly $1/\sigma$ of the node, key-rate
durations computed on the tension curve are concentrated near the bumped tenor: a five-year
risk shows up as five-year risk, not as a smear across the entire long end as it does under the
cubic. This yields hedge portfolios that are smaller, cheaper, and more stable, and it reduces
the cross-tenor leakage that inflates the apparent dimensionality of curve risk. The same
stability suppresses the spurious profit-and-loss and value-at-risk noise that interpolation
artefacts inject when a curve is re-struck daily. For no-arbitrage control, the tension dial
provides a direct lever to remove negative forwards where they arise. For model governance and
independent price verification, determinism, reproducibility, and a single documented parameter
are attractive in their own right. And the regulatory lineage through Smith--Wilson means the
energy formulation can be presented to supervisors in a language they already accept.

\subsection{When to use something else}
Intellectual honesty requires stating the method's limits. If all one needs is positive,
monotone forwards with minimal machinery, the monotone-convex method is simpler and is the
de-facto standard. If the goal is a parsimonious, economically interpretable description of the
curve --- level, slope, curvature --- rather than exact repricing, a Nelson--Siegel--Svensson
fit is the right tool. For Solvency~II discounting the Smith--Wilson form is prescribed,
although the present energy reading of its $\alpha$ may still aid intuition. Tension splines
also demand numerical care: the hyperbolic basis overflows at large $\sigma h$ and must be
evaluated in a stabilised form, and a single global tension, as the equipartition optimum here
illustrates, can over-tension a sparsely knotted long end while a densely knotted short end is
still curved --- the motivation for the per-segment calibration of Section~\ref{sec:nonunif}. Used within these
bounds, the method delivers a smooth, local, no-arbitrage curve under one transparent and
automatically chosen parameter.


\section{Discussion, Conclusion and Future Direction}\label{sec:discussion}
Placing the tension spline on its variational footing as a tensioned elastic beam achieves
three things. It \emph{derives} the governing equation \eqref{eq:euler} rather than asserting
it; it \emph{identifies} the smoothness and length norms as the flexural and axial energies of a
single elastic system, removing the conceptual double-counting of earlier beam-energy arguments;
and it \emph{motivates} a selection rule, \eqref{eq:phi}, that is dimensionless, range- and
scale-invariant, monotone (hence uniquely solvable), and equipped with a transparent
risk--hedging dial $\varphi^\star$. On the $22$-May-$2026$ U.S. Treasury curve the rule returns
a unique optimum ($\sigma^\star=4.86$ at equipartition) and, at that optimum, cuts the
hedge-locality radius of the benchmark $5$-year node from $6.7$ years to $0.8$ year while
removing the spurious long-end forward hump the cubic produces.

Several limitations and extensions remain. The global-tension construction of
Sections~\ref{sec:energy}--\ref{sec:validation} tensions the sparsely knotted long end as firmly
as the densely knotted short end; Section~\ref{sec:nonunif} resolves this with a per-segment
calibration that equalises the energy fraction across segments, and Section~\ref{sec:fwdnu}
combines it with the forward-curve construction of Section~\ref{sec:fwd}. We have
also worked within small-deflection beam theory, which is amply justified at the natural scale of
a rate curve, where $f'$ is minute; a geometrically exact (elastica) treatment, retaining the
arc-length element and the full curvature $f''/(1+(f')^2)^{3/2}$, would be required only if the
axes were rescaled to comparable magnitudes, and the scale-invariance of $\varphi$ makes that
unnecessary. Section~\ref{sec:multicurve} places the method inside a multi-curve (OIS/SOFR)
framework, with discounting bootstrapped exactly and the tension applied to the projection
curve; extending the consistency treatment to cross-currency and credit-sensitive bases is the
natural continuation. Modelling rate curves by analogy with
physical elastic systems remains, in our view, a fertile source of well-posed and
interpretable construction rules; the present paper is one such attempt.

\paragraph{Code and data availability.} A complete, reproducible reference implementation of every
method, figure, and result in this paper --- together with the historical-study and benchmark
tooling and the U.S. Treasury data used --- is available at
\url{https://github.com/ajsinha/fin-engg/tree/main/energy-partition-tension-spline}.
Running \texttt{python tension\_curve.py} regenerates the full validation report and figures; the
accompanying retrieval module reproduces the $2006$--$2026$ historical study from public Federal
Reserve (FRED) data.

\appendix
\section{Non-Uniform Tension Spline}\label{app:nonunif}
Allowing a distinct tension $\sigma_i$ on the interval $[x_i,x_{i+1}]$, the tension-spline ODE
$f^{(\mathrm{iv})}-\sigma_i^2 f''=0$ holds segment-wise, so $f$, $f'$ and $f''$ take the forms
\eqref{eq:tspline}--\eqref{eq:t2deriv} on each segment with $\sigma\!\to\!\sigma_i$ and
$h\!\to\!h_i$; in particular $f''(x)=[z_i\sinh(\sigma_i(x_{i+1}-x))+z_{i+1}\sinh(\sigma_i(x-x_i))]/\sinh(\sigma_i h_i)$,
so $f''$ is continuous at the knots by construction, with $f''(x_i)=z_i$. The unknowns
$\{z_i\}$ are fixed by continuity of $f'$ at each interior knot $x_i$. Specialising the slope
\eqref{eq:tderiv} to the left segment $[x_{i-1},x_i]$ (tension $\sigma_i$, width $h_{i-1}$) and
evaluating at the right end $x=x_i$ --- where $\cosh(\sigma_i(x_i-x))=\cosh 0=1$ and
$\cosh(\sigma_i(x-x_{i-1}))=\cosh(\sigma_i h_{i-1})$ --- gives
\[
f'(x_i^-)=\frac{-z_{i-1}+z_i\cosh(\sigma_i h_{i-1})}{\sigma_i\sinh(\sigma_i h_{i-1})}
+\frac{(y_i-y_{i-1})-(z_i-z_{i-1})/\sigma_i^2}{h_{i-1}} .
\]
Specialising instead to the right segment $[x_i,x_{i+1}]$ (tension $\sigma_{i+1}$, width $h_i$)
and evaluating at the left end $x=x_i$ --- where now $\cosh(\sigma_{i+1}(x_{i+1}-x))=\cosh(\sigma_{i+1}h_i)$
and $\cosh(\sigma_{i+1}(x-x_i))=1$ --- gives
\[
f'(x_i^+)=\frac{-z_i\cosh(\sigma_{i+1}h_i)+z_{i+1}}{\sigma_{i+1}\sinh(\sigma_{i+1}h_i)}
+\frac{(y_{i+1}-y_i)-(z_{i+1}-z_i)/\sigma_{i+1}^2}{h_i} .
\]
Set $f'(x_i^-)=f'(x_i^+)$ and collect the coefficient of each $z$, noting that the linear terms
$-(z_i-z_{i-1})/(\sigma_i^2 h_{i-1})$ and $-(z_{i+1}-z_i)/(\sigma_{i+1}^2 h_i)$ each contribute to
two coefficients. The $z_{i-1}$ terms (both from the left segment) combine as
$-1/(\sigma_i\sinh(\sigma_i h_{i-1}))+1/(\sigma_i^2 h_{i-1})$; the $z_{i+1}$ terms (both from the
right) as $1/(\sigma_{i+1}\sinh(\sigma_{i+1}h_i))-1/(\sigma_{i+1}^2 h_i)$ on the right side, which
moves to the left with a sign change; and the $z_i$ terms gather one contribution from each
segment. The $y$ terms move to the right-hand side as the centred second difference. The result
is the tridiagonal recursion, for $i=1,\dots,n-2$,
\begin{equation}
\mathcal{A}_{i-1}z_{i-1}+\mathcal{A}_i z_i+\mathcal{A}_{i+1}z_{i+1}
=\frac{y_{i+1}-y_i}{h_i}-\frac{y_i-y_{i-1}}{h_{i-1}},
\label{eq:nonunif}
\end{equation}
\begin{align}
\mathcal{A}_{i-1}&=\frac{1}{\sigma_i^2 h_{i-1}}-\frac{1}{\sigma_i\sinh(\sigma_i h_{i-1})},
\qquad
\mathcal{A}_{i+1}=\frac{1}{\sigma_{i+1}^2 h_i}-\frac{1}{\sigma_{i+1}\sinh(\sigma_{i+1}h_i)},\\
\mathcal{A}_i&=\left[\frac{\cosh(\sigma_i h_{i-1})}{\sigma_i\sinh(\sigma_i h_{i-1})}-\frac{1}{\sigma_i^2 h_{i-1}}\right]
+\left[\frac{\cosh(\sigma_{i+1}h_i)}{\sigma_{i+1}\sinh(\sigma_{i+1}h_i)}-\frac{1}{\sigma_{i+1}^2 h_i}\right].
\end{align}
The terms $1/(\sigma^2 h)$ arise from the linear part of the slope \eqref{eq:tderiv} and must
not be dropped. With the natural conditions $z_0=z_{n-1}=0$ the system is diagonally dominant
and solved exactly as in the uniform case. Setting all $\sigma_i=\sigma$ and $h_i=h$ gives
$\mathcal{A}_{i\pm1}=\sigma^{-2}\alpha$ and $\mathcal{A}_i=2\sigma^{-2}\beta$ with $\alpha,\beta$
as in \eqref{eq:ttridiag}, so the recursion reduces to $\sigma^{-2}$ times the uniform system
and returns the same $\{z_i\}$. The per-segment energies of Section~\ref{sec:nonunif} are then
evaluated from \eqref{eq:Eb} and \eqref{eq:Ea} with $\sigma\to\sigma_i$, and the calibration
\eqref{eq:segphi} is the fixed-point iteration described there.

\section{Closed-Form Energy Integrals}\label{app:energy}
Both per-segment energies follow from four elementary integrals. With $u\in[0,h]$ and the
substitutions below they are
\begin{align*}
\int_0^h\sinh^2(\sigma u)\,du&=\tfrac{\sinh(2\sigma h)}{4\sigma}-\tfrac h2, &
\int_0^h\cosh^2(\sigma u)\,du&=\tfrac{\sinh(2\sigma h)}{4\sigma}+\tfrac h2,\\
\int_0^h\sinh(\sigma(h-u))\sinh(\sigma u)\,du&=\tfrac12\!\left(h\cosh\sigma h-\tfrac{\sinh\sigma h}{\sigma}\right), &
\int_0^h\cosh(\sigma(h-u))\cosh(\sigma u)\,du&=\tfrac12\!\left(h\cosh\sigma h+\tfrac{\sinh\sigma h}{\sigma}\right),
\end{align*}
where the cross terms use $\sinh a\sinh b=\tfrac12[\cosh(a{+}b)-\cosh(a{-}b)]$ and
$\cosh a\cosh b=\tfrac12[\cosh(a{+}b)+\cosh(a{-}b)]$ with $a{+}b=\sigma h$ constant.

\emph{Flexural energy.} From \eqref{eq:t2deriv}, on $[x_i,x_{i+1}]$ with $u=x-x_i$,
$f''=\big[z_i\sinh(\sigma(h_i-u))+z_{i+1}\sinh(\sigma u)\big]/\sinh(\sigma h_i)$. Squaring and
integrating term by term with the first and third integrals above reproduces the per-segment
flexural energy \eqref{eq:Eb} quoted in Section~\ref{sec:energy}:
\[
E^{(i)}_{\mathrm{bend}}=\frac{1}{2\sinh^2(\sigma h_i)}
\left[(z_i^2+z_{i+1}^2)\!\left(\frac{\sinh(2\sigma h_i)}{4\sigma}-\frac{h_i}{2}\right)
+z_i z_{i+1}\!\left(h_i\cosh(\sigma h_i)-\frac{\sinh(\sigma h_i)}{\sigma}\right)\right].
\]
\emph{Axial energy.} From \eqref{eq:tderiv}, $f'=\big[-z_i\cosh(\sigma(h_i-u))+z_{i+1}\cosh(\sigma u)\big]/(\sigma\sinh\sigma h_i)+S_i$
with constant slope $S_i=\big[(y_{i+1}-y_i)-(z_{i+1}-z_i)/\sigma^2\big]/h_i$. Expanding $(f')^2$
into hyperbolic-squared, cross, and constant parts and using the second and fourth integrals,
together with $\int_0^{h_i}\!\big[-z_i\cosh(\sigma(h_i-u))+z_{i+1}\cosh(\sigma u)\big]du
=(z_{i+1}-z_i)\sinh(\sigma h_i)/\sigma$, gives
\begin{equation}
E^{(i)}_{\mathrm{tens}}
=\frac{1}{2\sinh^2(\sigma h_i)}
\left[(z_i^2+z_{i+1}^2)\!\left(\frac{\sinh(2\sigma h_i)}{4\sigma}+\frac{h_i}{2}\right)
-z_i z_{i+1}\!\left(h_i\cosh(\sigma h_i)+\frac{\sinh(\sigma h_i)}{\sigma}\right)\right]
+S_i\,(z_{i+1}-z_i)+\tfrac12\sigma^2 S_i^2 h_i .
\label{eq:Ea}
\end{equation}
Equations \eqref{eq:Eb} and \eqref{eq:Ea} give $\varphi(\sigma)=\sum_i E^{(i)}_{\mathrm{bend}}/\sum_i(E^{(i)}_{\mathrm{bend}}+E^{(i)}_{\mathrm{tens}})$
in closed form, with no quadrature required; the reference implementation verifies them against
$64$-point Gauss--Legendre quadrature to a relative error below $10^{-12}$.

\section{Numerically Stable Evaluation}\label{app:stable}
For large $\sigma h_i$ the factors $\sinh(\sigma h_i)$ and $\cosh(\sigma h_i)$ overflow, yet the
quantities that appear are bounded \emph{ratios}. Writing $u\in[0,h]$ and factoring $e^{\sigma h}$
from numerator and denominator,
\begin{equation}
\frac{\sinh(\sigma u)}{\sinh(\sigma h)}=\frac{e^{\sigma(u-h)}-e^{-\sigma(u+h)}}{1-e^{-2\sigma h}},
\qquad
\frac{\cosh(\sigma u)}{\sinh(\sigma h)}=\frac{e^{\sigma(u-h)}+e^{-\sigma(u+h)}}{1-e^{-2\sigma h}},
\end{equation}
in which every exponent is non-positive (since $0\le u\le h$), so no term exceeds unity and the
expressions are evaluated without overflow for arbitrarily large $\sigma h$. As $\sigma h\to\infty$
they tend to $e^{-\sigma(h-u)}$ and $e^{-\sigma(h-u)}$ respectively, recovering the
boundary-layer behaviour used in the proof of Proposition~\ref{prop:phi}. Every spline
evaluation in this paper uses these forms below a threshold $\sigma h<30$ and the direct
$\sinh,\cosh$ forms above it; the energy ratios in \eqref{eq:Eb}--\eqref{eq:Ea} are likewise
evaluated through $\coth(\sigma h)=1/\tanh(\sigma h)$ and $1/\sinh(\sigma h)$, both bounded.

\begin{thebibliography}{99}
\small
\bibitem{Anderson} L.~Andersen. \emph{Discount curve construction with tension splines.}
Review of Derivatives Research, 10(3):227--267, 2007.
\bibitem{Bianchetti} M.~Bianchetti. \emph{Two curves, one price.} Risk Magazine, August 2010.
\bibitem{Henrard} M.~Henrard. \emph{Interest Rate Modelling in the Multi-Curve Framework:
Foundations, Evolution and Implementation.} Applied Quantitative Finance. Palgrave Macmillan, 2014.
\bibitem{Ametrano} F.~M.~Ametrano, M.~Henrard. \emph{Everything you always wanted to know about
multiple interest rate curve bootstrapping but were afraid to ask.} Working paper, 2014.
\bibitem{Schweikert} D.~G.~Schweikert. \emph{An interpolation curve using a spline in tension.}
Journal of Mathematics and Physics, 45:312--317, 1966.
\bibitem{Cline} A.~K.~Cline. \emph{Scalar- and planar-valued curve fitting using splines under
tension.} Communications of the ACM, 17:218--223, 1974.
\bibitem{Costantini} P.~Costantini, B.~I.~Kvasov, C.~Manni. \emph{On discrete hyperbolic tension
splines.} Advances in Computational Mathematics, 11(4):331--354, 1999.
\bibitem{Hagan} P.~S.~Hagan, G.~West. \emph{Interpolation methods for curve construction.}
Applied Mathematical Finance, 13(2):89--129, 2006.
\bibitem{Renka} R.~J.~Renka. \emph{Interpolatory tension splines with automatic selection of
tension factors.} SIAM J.\ Sci.\ Stat.\ Comput., 8(3):393--415, 1987.
\bibitem{Kamada} M.~Kamada, R.~Enkhbat. \emph{Spline interpolation in piecewise constant
tension.} SAMPTA'09, Marseille, 2009.
\bibitem{McCulloch} J.~H.~McCulloch. \emph{Measuring the term structure of interest rates.}
Journal of Business, 44:19--31, 1971.
\bibitem{Tanggaard} C.~Tanggaard. \emph{Nonparametric smoothing of yield curves.}
Review of Quantitative Finance and Accounting, 9:251--267, 1997.
\bibitem{Fritsch} F.~N.~Fritsch, R.~E.~Carlson. \emph{Monotone piecewise cubic interpolation.}
SIAM J.\ Numer.\ Anal., 17(2):238--246, 1980.
\bibitem{Ron} U.~Ron. \emph{A practical guide to swap curve construction.} Bank of Canada
Working Paper 2000-17.
\bibitem{Kincaid} D.~Kincaid, W.~Cheney. \emph{Numerical Analysis: Mathematics of Scientific
Computing.} American Mathematical Society.
\bibitem{TreasuryMethod} U.S. Department of the Treasury. \emph{Treasury Par Yield Curve
Methodology}, 2021 (monotone convex spline, effective 6-Dec-2021).
\bibitem{H15} Board of Governors of the Federal Reserve System. \emph{Selected Interest Rates
(H.15)}, daily release, $22$-May-$2026$.
\end{thebibliography}

\end{document}
